\documentclass[11pt]{article}

\usepackage[T1]{fontenc}
\usepackage[utf8]{inputenc}
\usepackage[english]{babel}
\usepackage[margin=1in]{geometry}
\usepackage{lmodern}
\usepackage{microtype}
\usepackage{mathtools}
\usepackage{amsmath,amssymb,amsthm}
\usepackage{array}
\usepackage{booktabs}
\usepackage{tabularx}
\usepackage{longtable}
\usepackage{float}
\usepackage{graphicx}
\usepackage{tikz}
\usepackage{pdflscape}
\usepackage{caption}
\usepackage{enumitem}
\usepackage{fancyhdr}
\usepackage[numbers]{natbib}
\usepackage{doi}
\usepackage{xurl}
\usepackage{aliascnt}
\usepackage{hyperref}
\usepackage{cleveref}
\usepackage{orcidlink}
\usepackage{titling}

% Meta-math paper macros.

%--- Inherited substrate/foundations bibkey markers --------------------------
\newcommand{\FOne}{Foundations~I}
\newcommand{\FTwo}{Foundations~II}
\newcommand{\FThree}{Foundations~III}

%--- Inherited apparatus (cited; not coined here) ----------------------------

%--- Part I (Foreclosure) state predicates ----------------------------------
\newcommand{\Xequiv}{$X$-equivalent}
\newcommand{\nonXequiv}{non-$X$-equivalent}
\newcommand{\CRETE}{\textup{\textsc{CRE-TE}}}
\newcommand{\CREBF}{\textup{\textsc{CRE-BF}}}
\newcommand{\CRECTMT}{\textup{\textsc{CRE-CTMT}}}
\newcommand{\XstrengthObl}{$X$-strength obligation}

%--- Part I framework-internal moves ----------------------------------------
\newcommand{\fim}{framework-internal move}
\newcommand{\carrierPivot}{carrier pivot}
\newcommand{\bridgeImport}{bridge import}
\newcommand{\cascadeInternal}{cascade-internal computation}
\newcommand{\TCD}{typed-context discipline}
\newcommand{\TypedBridgeCarrier}{carrier-side typed bridge}

%--- Part II V-Differential column markers ----------------------------------
\newcommand{\VDTSO}{V-Differential trace-state-only column}
\newcommand{\VDWCE}{V-Differential witness-content-exposing column}

%--- Part II three-option derivation sweep + verdict signature --------------
\newcommand{\SigPattern}{$\mathrm{(ii)}/\mathrm{(iii)}/\mathrm{(ii)}$}

%--- Part II recognition vs derivation mode ---------------------------------
\newcommand{\RML}{recognition-mode landing}

%--- Part III calibration seed catalog --------------------------------------
\newcommand{\Cvobs}{\mathcal{V}_{\text{obs}}}
\newcommand{\CMI}{\mathcal{M}_{\text{I}}}
\newcommand{\CRE}{\mathcal{R}_{\text{E}}}
\newcommand{\CCh}{\mathcal{C}_{\text{h}}}
\newcommand{\CQamp}{\mathcal{Q}_{\text{amp}}}
\newcommand{\CIalg}{\mathcal{I}_{\text{alg}}}
\newcommand{\CRreal}{\mathcal{R}_{\text{real}}}
\newcommand{\COCstar}{\mathcal{O}_{C^*}}
\newcommand{\CFexistsSO}{\mathcal{F}_{\exists \mathrm{SO}}}
\newcommand{\CSHO}{\mathcal{S}_{\mathrm{HO}}}
\newcommand{\CPPH}{\mathcal{P}_{\mathrm{PH}}}
\newcommand{\CAVNP}{\mathcal{A}_{\mathrm{VNP}}}
\newcommand{\CMmot}{\mathcal{M}_{\text{mot}}}

%--- Part III calibration apparatus ----------------------------------------
\newcommand{\fourAxis}{four-axis calibration profile}
\newcommand{\TypedHole}{typed coverage hole}
\newcommand{\xiNCone}{h_{\mathrm{NC}^{1}}}
\newcommand{\xiAdaptive}{h_{\mathrm{adapt}}}
\newcommand{\FType}{framework-type}
\newcommand{\TDS}{typed-deposit species}
\newcommand{\TypedBridgeCalibration}{calibration-side typed bridge}

%--- Part III named functorial gates ---------------------------------------
\newcommand{\GateFunctionDegree}{G_{\mathrm{deg}}}
\newcommand{\GatePlantedStatistical}{G_{\mathrm{plant}}}
\newcommand{\GateUniformTime}{G_{\mathrm{unif}}}

%--- Lean-disclosure markers (body uses only obligation-ref; appendix D uses LeanDecl) ---
\newcommand{\LeanDecl}[1]{\texttt{\detokenize{#1}}}

% Semantic brackets for the interpretation of claims.
\providecommand{\llbracket}{[\![}
\providecommand{\rrbracket}{]\!]}


\newtheorem{theorem}{Theorem}[section]
% Shared numbering through alias counters, so that cleveref names each
% environment correctly.
\newcommand{\sharedtheorem}[2]{%
  \newaliascnt{#1}{theorem}\newtheorem{#1}[#1]{#2}\aliascntresetthe{#1}}
\sharedtheorem{lemma}{Lemma}
\sharedtheorem{proposition}{Proposition}
\sharedtheorem{corollary}{Corollary}
\sharedtheorem{schema}{Conditional schema}
\theoremstyle{definition}
\sharedtheorem{definition}{Definition}
\sharedtheorem{remark}{Remark}
\floatstyle{ruled}
\newfloat{algorithm}{tbp}{loa}
\floatname{algorithm}{Algorithm}

\urlstyle{same}
% Lean identifiers: break only after dots and underscores.
\DeclareRobustCommand{\leanid}[1]{{\def\UrlBreaks{\do\.\do\_}\def\UrlBigBreaks{}%
  \def\UrlNoBreaks{}\nolinkurl{#1}}}
\captionsetup{
  font=small,
  labelfont=bf,
  labelsep=period
}
\captionsetup[table]{skip=6pt,position=top}
\captionsetup[figure]{skip=6pt}
\captionsetup[algorithm]{skip=6pt,position=top}
\crefname{schema}{conditional schema}{conditional schemas}
\Crefname{schema}{Conditional schema}{Conditional schemas}
\crefname{algorithm}{algorithm}{algorithms}
\Crefname{algorithm}{Algorithm}{Algorithms}
\crefname{theorem}{theorem}{theorems}
\Crefname{theorem}{Theorem}{Theorems}
\crefname{lemma}{lemma}{lemmas}
\Crefname{lemma}{Lemma}{Lemmas}
\crefname{proposition}{proposition}{propositions}
\Crefname{proposition}{Proposition}{Propositions}
\crefname{corollary}{corollary}{corollaries}
\Crefname{corollary}{Corollary}{Corollaries}
\crefname{definition}{definition}{definitions}
\Crefname{definition}{Definition}{Definitions}
\crefname{remark}{remark}{remarks}
\Crefname{remark}{Remark}{Remarks}
\setlist[enumerate]{leftmargin=*, itemsep=2pt, topsep=4pt}
\setlength{\droptitle}{-3.2em}
\pretitle{\begin{center}\LARGE\bfseries}
\posttitle{\par\end{center}\vspace{0.5em}}
\preauthor{\begin{center}\normalsize}
\postauthor{\par\end{center}\vspace{-0.4em}}
\predate{\begin{center}\small}
\postdate{\par\end{center}\vspace{0.6em}}
\setlength{\emergencystretch}{2em}
\clubpenalty=10000
\widowpenalty=10000
\displaywidowpenalty=10000
\interfootnotelinepenalty=10000
\allowdisplaybreaks[2]
\fancypagestyle{firstpage}{\fancyhf{}
  \fancyfoot[C]{\scriptsize Automorph Inc., Wilmington, DE, USA \quad \texttt{ioannis@automorph.io} \quad \textcopyright\ 2026 \quad CC-BY 4.0 \quad \doi{10.5281/zenodo.23086999}}
  \renewcommand{\headrulewidth}{0pt}
  \renewcommand{\footrulewidth}{0pt}
}

\title{Three Conditional Clay-Problem Closures: A Shared Gaussian Observation and a Common Vocabulary}
\author{Ioannis Tsiokos\,\orcidlink{0009-0009-7659-5964}}
\date{Version 4: October 1, 2026\\[2pt] \footnotesize (Version 1: June 17, 2026)}
\hypersetup{
  pdftitle={Three Conditional Clay-Problem Closures: A Shared Gaussian Observation and a Common Vocabulary},
  pdfauthor={Ioannis Tsiokos},
  hidelinks
}

\begin{document}

\maketitle
\thispagestyle{firstpage}

\begin{abstract}
Three companion papers state conditional closures of Navier--Stokes
regularity, the Riemann hypothesis and \(P\) versus \(NP\) within Six
Birds Theory, each under its own open premises.  This paper compares
them in two ways.  Concretely, one family of complex Gaussian
observations can be evaluated on all three settings.  For positive
width, viscosity and time, and for a finite window of zeta zeros closed
under reflection in the critical line with reflection-invariant
nonnegative integer weights, an elementary computation shows that the heat excess of the
window is a positive multiple of its residual charge, the weighted sum
of squared distances of the zeros from the critical line.  Applied to a
mild Navier--Stokes solution, the observation restates the mild
equation with its nonlinear Duhamel term; on the assignments of a CNF
formula the same Gaussian kernel gives a mass that is positive exactly
for satisfiable formulas, as it would for any \(0/1\)-valued filter.  These identities are formally verified in
Lean.  They share a formula, not a mechanism, and transfer no premise
from one problem to another.  Structurally, the paper gives a
vocabulary for proof attempts and closures: a four-state classification
of attempts aimed at a target, four conditional schemas for closures
that name their load-bearing source instead of deriving it, and a
catalogue of typed obstruction detectors with three named open
transports.  Its numbered results in these parts are logical
consequences of explicitly assumed schemas or checks of declared finite
data.  No unconditional result about any of the problems is claimed.
\end{abstract}

\medskip
\noindent\textbf{Keywords.} Six Birds Theory; conditional closure;
Riemann hypothesis; Navier--Stokes equations; P versus NP; Gaussian
observation; adequacy residual; proof attempts; Lean formalization.

\section{Introduction}\label{sec:intro}

Six Birds Theory is a program that treats a theory as a fixed point of
an idempotent closure operator acting on descriptions of a system, and
studies how such closed descriptions arise
\cite{Tsiokos2026EmergenceCalculus}.  Its companion papers apply this
vocabulary to open problems (a glossary is in
\Cref{sec:glossary}): a \emph{conditional closure} of
a problem is a theorem deriving the problem's standard statement from
explicitly named premises.  Three companion papers state conditional
closures of three Clay Millennium problems in this sense
\cite{Fefferman2000NavierStokes,Bombieri2000Riemann,Cook2000PvsNP,CarlsonJaffeWiles2006MillenniumPrize}.
The Navier--Stokes paper \cite{Tsiokos2026NavierStokes} derives a
global solution satisfying the equations pointwise from a spectral
closure hypothesis and hypotheses on the initial data.  The Riemann-hypothesis paper
\cite{Tsiokos2026RH} derives the Riemann hypothesis from a supplied
Selberg recognition source and an identification with the classical
zeros.  The P-versus-NP paper \cite{Tsiokos2026PvNP} proves a fixed
machine-level translation between deciding SAT and computing a tagged
witness selector, and derives \(P\neq NP\) from an accepted negative
source.  In each case the premises are open.  None of the three is an
unconditional result, and this paper does not change that.  The three
are also not of one type: the Navier--Stokes closure is derived from its
hypotheses through problem-specific analysis, while the other two name
a load-bearing source as a premise.

The three arguments nevertheless have a common shape, and this paper
asks why.  It gives two kinds of answer.  The first is concrete: one
family of Gaussian observations can be read on all three mathematical
objects (\Cref{sec:gaussian-bridge,sec:triad-bridge}).  The second is
structural: a description, in the vocabulary of the framework, of how
proof attempts aimed at a target can stall, of how a closure can instead
name its load-bearing source, and of the detectors by which obstructions
are measured (\Cref{sec:foreclosure,sec:recognition,sec:calibration}).

\paragraph{The concrete construction.}  For a width \(a>0\) and a
complex centre \(z\), a complex Gaussian window on the Fourier side of
\(\mathbb R^{3}\) defines a linear observation of Fourier states.
Evaluated against the heat multiplier, it gives
\((\pi/(a+b))^{3/2}e^{-\kappa z^{2}}\) with \(b=4\pi^{2}\nu t\) and
\(\kappa=ab/(a+b)\).  Placing \(z\) at the spectral coordinate of a zero
\(\rho\) of the completed zeta function, the logarithmic excess of this
value is \(\kappa(\operatorname{Re}\rho-\frac12)^{2}\); summed with
reflection-invariant nonnegative integer weights over a finite window closed under
reflection in the critical line, it is \(\kappa\) times
the residual charge of the window in the sense of the
adequacy-residual companion \cite{Tsiokos2026XiCompanion}
(\Cref{thm:joint-gaussian-return}).  Vanishing of this excess on all
windows is a reformulation of the Riemann hypothesis, not a route to
it.  Applied to a mild solution of the Navier--Stokes equations the
same observation restates the mild equation, with the nonlinear Duhamel
term kept and not estimated.  On the assignment space of a CNF formula
the same scalar kernel gives a Gaussian mass that is positive exactly
for satisfiable formulas, with exact laws for adding a clause and for
fixing a variable (\Cref{thm:concrete-triad-gaussian-xi}); positivity
here holds for any \(\{0,1\}\)-valued filter and is not specific to
SAT.  On a reflected
zero window the constant probe reads zero; on a satisfiable formula it
reads the number of satisfying assignments.  These identities are
elementary, and they are formally verified in Lean
(\Cref{app:formalization}).  What the three settings share is a
formula, not a mechanism: the identities transfer no premise from one
problem to another, and no efficient evaluation of the Gaussian mass is
proved.

\paragraph{The structural account.}  The account has three parts.
Part~I (\Cref{sec:foreclosure}) concerns \emph{carriers}: a carrier is a
presentation through which a proof attempt tries to reach a target
\(X\), together with the obstruction it leaves, its \emph{residual}.
Four outcomes are distinguished: the residual closure points to a
proposition other than \(X\); it is equivalent to \(X\) but the attempt
can certify \(X\) only by auditing itself; it needs a bridge to \(X\)
that is unavailable; or its reduction ends at a family of matrix
elements whose decision it does not supply.  That these four outcomes
are exhaustive is the Coverage schema, which is not proved; the
theorems of Part~I derive consequences from explicitly assumed schemas
and are short logical arguments.

Part~II (\Cref{sec:recognition}) concerns closures that do not derive
their load-bearing content but name it as a premise, called
\emph{recognition-mode landings}.  It states, as four conditional
schemas, a seven-stage template for such closures, a pattern in the
outcome of searching for an internal derivation, the thesis that the
load-bearing content is named rather than derived, and a combination of
this with the foreclosure statement of Part~I.  The Riemann-hypothesis
and P-versus-NP applications illustrate the schemas; they do not prove
them.

Part~III (\Cref{sec:calibration}) concerns \emph{calibrations}, typed
detectors of obstructions such as low-degree polynomials, mutual
information or Nullstellensatz certificates.  It records a catalogue,
introduced here, of thirteen detector classes considered for the
P-versus-NP problem, a criterion for
when a new calibration is a relabelling of an old one, the direction of
bridges between calibrations, a table of which regions the catalogue
covers, and three named transports---\emph{functorial gates}---whose
proofs are open.  The results of Part~III are elementary facts about
these finite records; they say nothing about the performance of a
detector or about inclusions between complexity classes.

\paragraph{Status of the results.}  The paper distinguishes three
kinds of numbered statement (\Cref{sec:apparatus:opacity-guards}).
Theorems, propositions and corollaries are proved from the hypotheses
they state; in Parts~I and~III most of them are short logical arguments
or checks of declared finite data, and we say so where they occur.  Conditional schemas are not proved; they are displayed so
that other statements can assume them.  Some predicates, such as being
in framework scope or being load-bearing, are primitive: they are named
and not defined here, and results about them establish only their
logical shape.  \Cref{app:formalization} records, for every numbered
statement, its formal counterpart in Lean.  The only computations in
the paper are in \Cref{sec:gaussian-bridge,sec:triad-bridge}, and they
are elementary; their contribution is that the definitions are aligned
across the three settings and formally verified.  The rest of the paper
is a vocabulary for comparing conditional closures, and a reader may
take it as such.
A glossary of the framework's terms is in
\Cref{sec:glossary}.  The companion papers are preprints.

\paragraph{Relation to the series.}  The paper uses the closure
calculus of \FOne\ \cite{Tsiokos2026EmergenceCalculus}, the
admissibility setting of \FTwo\ \cite{Tsiokos2026FoundationsII}, and the
audit calculus of \FThree\ \cite{Tsiokos2026FoundationsIII}.  The three
applications \cite{Tsiokos2026NavierStokes,Tsiokos2026RH,Tsiokos2026PvNP}
are cited as they stand; the recognition sources of the two
trace-state-only applications are named in those applications.  The
self-dual trace confinement companion \cite{Tsiokos2026SDTC} supplies
the confinement squeeze and its setting, the hiddenness companion
\cite{Tsiokos2026CSL} the placements and exposure classifications used
here, the adequacy-residual companion \cite{Tsiokos2026XiCompanion} supplies
the residual, and the emergence companion \cite{Tsiokos2026NeedleKiller}
supplies the distinction between native and layer-dissolving probes.
This paper is downstream of all of these and replaces none of their
arguments.

\paragraph{Scope.}  The Birch--Swinnerton-Dyer conjecture
\cite{Wiles2000BSD} is outside the scope of the paper, and the Hodge
conjecture \cite{Deligne2000Hodge} appears only as a prediction.  Nothing
here asserts an unconditional proof of any Millennium problem
(\Cref{sec:scope}).

\paragraph{Organization.}  \Cref{sec:apparatus} fixes the vocabulary.
\Cref{sec:foreclosure,sec:recognition} contain Parts~I and~II, and
\Cref{sec:validation} places the three applications in that scheme.
\Cref{sec:gaussian-bridge,sec:triad-bridge} contain the Gaussian
construction, and \Cref{sec:calibration} Part~III.
\Cref{sec:audit} checks the description of the companion papers,
\Cref{sec:predictions} lists predictions and open problems, and
\Cref{sec:scope} collects the limitations.
\Cref{app:formalization} describes the formal verification and
\Cref{app:version-history} the changes between versions.

\section{Apparatus}\label{sec:apparatus}

This section fixes the vocabulary used in the rest of the paper.  Most
of it is imported from earlier papers of the series and is used as a
black box; we state only what the later sections need.

\subsection{Closure formation}\label{sec:apparatus:closure-formation}

Six Birds Theory treats a theory as a fixed point of an idempotent
closure operator acting on descriptions; this calculus is developed in
\FOne\ \cite{Tsiokos2026EmergenceCalculus}.  The companion papers apply
it to a mathematical problem by recording, for a \emph{substrate} (the
problem's mathematical setting), the carriers, observables, comparison
maps and audit data visible at a chosen layer.  When this typed record
is assumed to be admissibly closed under the framework's formation
discipline, we call it a \emph{formed closure}, written
\(\operatorname{FormedClosure}(S)\); the companion papers use the term
informally, and the predicate is a definition of this paper.
We call the standing assumption that the formed closures used by the
companion papers exist the \emph{SB closure assumption}.  It is a
framework assumption, not a theorem of ordinary set theory; it is not
proved here, and by itself it supplies none of the problem-specific
hypotheses of the three applications.

An \emph{accepted import} is outside mathematical content that the
closure formation allows a proof attempt to use.  It must be named and
kept apart from what the attempt produces itself.  To keep names
distinct from truth values, \Cref{sec:foreclosure} works with a set
\(\mathsf{Claim}\) of \emph{claims}, each with an interpretation
\(\llbracket c\rrbracket\) as a proposition.  The accepted imports form
a set of claims with a soundness hypothesis,
\[
  \mathcal B\subseteq\mathsf{Claim},\qquad
  c\in\mathcal B\ \Rightarrow\ \llbracket c\rrbracket .
\]
Two different claims may have equivalent interpretations, and a true
claim need not be accepted.  (A predicate on propositions would not
allow this: since equivalent propositions are equal, a sound predicate
on propositions accepts either every true proposition or none.)  The phrase
``imported as named recognition'', used for the recognition sources of
\Cref{sec:recognition}, means that a source is admitted in this way and
is not claimed to be derived.

\subsection{The Foundations III audit calculus}\label{sec:apparatus:birdint}

\FThree\ \cite{Tsiokos2026FoundationsIII} constructs a finite audited
calculus, \(\mathbf{BirdInt}^{\mathrm{aud}}_{\mathrm{fin}}\), in which
claims are recorded together with the instrument that audits them and a
verdict.  We use one of its results as motivation.  Theorem~24 of
\FThree\ (same-level self-audit failure) concerns the claim classifier
of that calculus: a well-formed complete self-soundness claim of an
instrument, judged at its own level, is not accepted; it receives
\texttt{outside\_scope}, \texttt{undefined\_circular}, or a status of
higher priority (\texttt{blocked}, \texttt{absent\_with\_record} or
\texttt{failed\_audit}).  In \Cref{sec:foreclosure} each proof attempt
\(C\) records the verdict of its same-level audit, and
\(\operatorname{SelfAuditFail}(C)\) is the statement that this recorded
verdict is \texttt{outside\_scope} or \texttt{undefined\_circular}; the
restriction to these two verdicts is an assumption of this paper, not a
consequence of Theorem~24.  The verdict
is part of the data of \(C\); it is not derived here from Theorem~24,
which explains why it is the relevant one: an attempt that can certify
its target only by auditing itself at the target's level has fallen
back on the content it was meant to supply.  The admissibility
background for these audits is that of \FTwo\
\cite{Tsiokos2026FoundationsII}.

\subsection{Adequacy residuals}\label{sec:apparatus:adequacy-residual}

The adequacy-residual companion \cite{Tsiokos2026XiCompanion} measures
how much of a target readout is left unexplained by a family of native
probes.  A finite-dimensional real or complex carrier space \(E\)
carries an \emph{audit energy} \(A\succeq0\), the cost of \(u\in E\)
being \(u^{*}Au\); the native probes form an operator \(L\) on \(E\), and
the target readout an operator \(D\), both required to vanish on
\(\ker A\) (\emph{null legality}: \(\ker A\subseteq\ker L\cap\ker D\)).  With the block currencies \(K_{FG}=FA^{\dagger}G^{*}\),
the residual is the generalized Schur complement
\[
  \Xi_{A}(D\mid L)=K_{DD}-K_{DL}K_{LL}^{\dagger}K_{LD}.
\]
The companion writes the energy as a subscript \(C\); here \(C\) is
reserved for carriers.  When the target readout is the identity and the
audit energy is the identity, \(\Xi(L):=\Xi_{I}(I\mid L)\) is the
orthogonal projection onto \(\ker L\), and
\Cref{sec:gaussian-bridge,sec:triad-bridge} use it in that form.  In
\Cref{sec:foreclosure} only one fact about the residual is used, namely
whether it vanishes.

\begin{definition}[Carrier residual]\label{def:carrier-residual}
For a carrier with audit energy \(A\), native probes \(L\) and target
readout \(D\), the \emph{carrier residual} is the adequacy residual
\(\Xi_{A}(D\mid L)\) of \cite{Tsiokos2026XiCompanion}; under null
legality it is the part of the target readout that the native probes do
not explain at the audit energy \(A\).
\end{definition}

A \emph{residual closure} is the statement \(\Xi_{A}(D\mid L)=0\).
When \Cref{sec:foreclosure} compares a residual closure with a target,
the comparison is logical: it asks whether the proposition
\(\Xi_{A}(D\mid L)=0\) is equivalent to the target.  A residual closure
does not by itself prove the target.

From the emergence companion \cite{Tsiokos2026NeedleKiller} we use the
distinction between probes that stay native to a carrier and probes that
\emph{dissolve the layer} in which the carrier operates.  A
layer-dissolving probe is not evidence that the original carrier closed
its target; its endpoint must either be recorded as outside content or
be admitted by the closure formation.

\subsection{Typed contexts}\label{sec:apparatus:typed-context}

Every claim, carrier, bridge, calibration and obligation in this paper
carries an explicit type: its source, its target, the kind of carrier
it uses, and its admissible imports.  This \emph{\TCD}, inherited from
\FTwo\ and \FThree\
\cite{Tsiokos2026FoundationsII,Tsiokos2026FoundationsIII}, forbids
replacing one of these by an informally similar one.  We write
\[
  C\leadsto X,\qquad C\models\text{\TCD}
\]
for ``the carrier \(C\) is aimed at the target \(X\)'' and ``\(C\) is
audited in a context whose source, target, carrier type and imports are
fixed''.  Two consequences are used repeatedly: a carrier aimed at one
target is not a carrier aimed at another without an explicit bridge or
re-aiming step, and a bridge from one setting to another does not by
itself give a bridge back.

\subsection{Theorems and conditional schemas}\label{sec:apparatus:opacity-guards}

The numbered statements of this paper are of two kinds.  A
\emph{theorem}, proposition or corollary is proved, from the hypotheses
it states.  A \emph{conditional schema} is a statement that is not
proved; it is displayed in its own environment, with its hypotheses, so
that later results can name it as a hypothesis.  Several predicates in
\Cref{sec:foreclosure,sec:recognition} are \emph{primitive}: they are
named but not defined in this paper, and nothing is assumed about them
beyond what a statement says explicitly.  A result about primitive
predicates establishes only the logical shape it displays.
\Cref{app:formalization} lists, for each numbered statement, what has
been formally verified.

\subsection{Glossary}\label{sec:glossary}

The following framework terms recur.  Each is used only in the sense
given here.

\begin{description}[leftmargin=1.5em,style=nextline,itemsep=2pt]
\item[Substrate.] The mathematical setting of a problem, for instance
  the zeros of the completed zeta function or the runs of machines on
  SAT instances.
\item[Formed closure; SB closure assumption.] A typed record of a
  substrate's carriers, observables, comparison maps and audit data,
  assumed admissibly closed (\Cref{sec:apparatus:closure-formation}); the
  assumption that the formed closures used by the companion papers
  exist.
\item[Carrier; residual.] A presentation through which a proof attempt
  tries to reach a target; the part of the target readout that the
  carrier's native probes leave unexplained
  (\Cref{def:carrier-residual}).
\item[Accepted import; claim.] Outside content that a closure admits; it
  is recorded as a claim, a name with an interpretation as a
  proposition (\Cref{sec:apparatus:closure-formation}).
\item[Layer-dissolving probe.] A probe whose endpoint lies outside the
  layer in which the carrier operates, so that it is not evidence that
  the carrier closed its target.
\item[Recognition source; recognition-mode landing.] A load-bearing
  structural fact that a closure names as a premise rather than
  deriving; a conditional closure built on such a premise
  (\Cref{sec:recognition}).
\item[Trace-state-only; witness-content-exposing.] The two columns of
  \Cref{def:vdiff-tso-column}: whether current observation sees only
  trace bookkeeping of a computation, or a witness datum itself.  For
  \(P\) versus \(NP\), current observation by polynomial-time machines
  is placed in the first column: it is not assumed to see a satisfying
  assignment.
\item[Calibration; functorial gate.] A typed detector of obstructions
  together with its native carrier; a named transport between
  calibrations or types whose proof is open (\Cref{sec:calibration}).
\end{description}

\section{Part I --- Foreclosure}\label{sec:foreclosure}

\subsection{Setup}\label{sec:foreclosure:setup}

Part~I describes how a proof attempt aimed at a target can stall once
the target, the attempt and its admissible imports have been fixed.
The target is a claim \(x\in\mathsf{Claim}\) in the sense of
\Cref{sec:apparatus:closure-formation}, with interpretation
\(\llbracket x\rrbracket\); the accepted imports form a sound set
\(\mathcal B\subseteq\mathsf{Claim}\).

A proof attempt is represented by a \emph{carrier}: a presentation
through which the attempt tries to reach \(x\), together with the
residual obstruction it leaves.  The carrier also records a few
target-relative facts about itself, such as whether an accepted bridge
to the target exists.  These facts are data of the carrier, not
consequences derived here; the results of this part show what follows
from them and from explicitly stated schemas.  They are not a claim
that every possible argument for \(x\) has this form.

\begin{definition}[Typed primary carrier]\label{def:typed-primary-carrier}
A \emph{typed primary carrier} \(C\) consists of
\begin{enumerate}
\item rational matrices: an audit energy \(A_{C}\), a native-probe
  matrix \(L_{C}\), a target (dissolving) readout \(D_{C}\) and a
  declared matrix \(H_{C}\) standing for \(K_{LL}^{\dagger}\), which
  determine the carrier residual
  \[
    R_{C}=K_{DD}-K_{DL}H_{C}K_{LD},\qquad K_{FG}=FA_{C}^{\natural}G^{\mathsf T},
  \]
  where \(A_{C}^{\natural}\) inverts the nonzero diagonal entries of
  \(A_{C}\) (this is the residual of \Cref{def:carrier-residual} when
  \(A_{C}\) is diagonal and \(H_{C}=K_{LL}^{\dagger}\));
\item the verdict of its same-level audit
  (\Cref{sec:apparatus:birdint});
\item target-relative predicates on claims: \(C\leadsto x\) (aimed
  at \(x\)), \(\operatorname{NoBridge}(C,x)\) (no accepted bridge to
  the target), \(\operatorname{BridgeMismatch}(C,x)\) (structural
  incompatibility with the target), and
  \(\operatorname{MDEqTarget}(C,x)\) (the decision of its terminal
  matrix family is equivalent to the target);
\item target-independent cascade data: \(\operatorname{MRed}(C)\) (its
  cascade reduces to a family of matrix elements) and
  \(\operatorname{IRDec}(C)\) (the inherited records decide that family);
\end{enumerate}
and the proposition \(C\models\text{\TCD}\).  The \emph{residual
closure} of \(C\) is the proposition \(R_{C}=0\).
\end{definition}

The residual is used only through the proposition \(R_{C}=0\); the
matrix data serve to make it a definite proposition.  Its
interpretation as an unexplained currency in the sense of
\Cref{sec:apparatus:adequacy-residual} needs the further hypotheses
just stated, which Part~I does not use, and \(R_{C}\) is a different
object from the Hilbert-space residual \(\Xi(L)\) of
\Cref{sec:gaussian-bridge,sec:triad-bridge}.

\begin{definition}[Four typed carrier states]\label{def:four-typed-states}
Let \(x\) be a target claim and \(C\) a carrier.  \(C\) is
\emph{\Xequiv} if \((R_{C}=0)\leftrightarrow\llbracket x\rrbracket\),
and \emph{\nonXequiv} if its residual closure is equivalent to a
different claim:
\[
  \mbox{\nonXequiv}_{x}(C)
  \;\Longleftrightarrow\;
  \exists c'\in\mathsf{Claim}.\;
    \bigl((R_{C}=0)\leftrightarrow\llbracket c'\rrbracket\bigr)
    \wedge c'\neq x .
\]
The three \emph{carrier-residual-equivalent} (CRE) states, named
\emph{target-equivalent} (TE), \emph{bridge-failed} (BF) and
\emph{cascade-terminal matrix-element} (CTMT), are
\[
\begin{aligned}
  \mbox{\CRETE}_{x}(C)
    &\Longleftrightarrow
      \mbox{\Xequiv}_{x}(C)\wedge\operatorname{SelfAuditFail}(C),\\
  \mbox{\CREBF}_{x}(C)
    &\Longleftrightarrow
      \operatorname{NoBridge}(C,x)\vee\operatorname{BridgeMismatch}(C,x),\\
  \mbox{\CRECTMT}_{x}(C)
    &\Longleftrightarrow
      \operatorname{MRed}(C)\wedge\neg\operatorname{IRDec}(C)
      \wedge\neg\operatorname{MDEqTarget}(C,x):
\end{aligned}
\]
that is, the attempt can certify the target only by a failed
self-audit; it has no usable bridge to the target; or its reduction ends
at a family of matrix elements that its inherited records do not
decide.  For a claim \(c'\), put
\[
  \operatorname{NativeOutcome}_{\mathcal B}(C,c')
  \;\Longleftrightarrow\;
  \bigl(c'\in\mathcal B\wedge R_{C}=0\bigr)\vee c'\notin\mathcal B ,
\]
so that either \(c'\) is accepted and the carrier closes natively, or
\(c'\) is not accepted; we call a claim that is not accepted
\emph{open-pending}, whatever its truth value.
\end{definition}

Because claims are compared as claims, a carrier can be both
\nonXequiv\ and \Xequiv: its residual closure may be equivalent to the
target and to a different claim with the same truth value.  In a claim
language rich enough to contain, with every claim, a differently named
equivalent one, every \Xequiv\ carrier is also \nonXequiv.  The
intended reading takes claims to be statements up to a fixed
normalization, but no result below depends on this choice, and the
substance of the classification lies in the hypothesis (Cls) below
rather than in the first state.

\begin{definition}[Coverage context]\label{def:coverage-ctx}
\(\operatorname{CoverageCtx}(x,C)\) is a primitive predicate marking the
audit settings in which the four states are intended to be exhaustive.
\end{definition}

\subsection{Carrier Dichotomy}\label{sec:foreclosure:carrier-dichotomy}

Write \(\operatorname{ExactlyOne}(A,B,D)\) for
\((A\wedge\neg B\wedge\neg D)\vee(B\wedge\neg A\wedge\neg D)\vee(D\wedge\neg A\wedge\neg B)\).
The \emph{classification hypothesis} for \((x,C)\) is
\[
  \mbox{(Cls)}\qquad
  \mbox{\Xequiv}_{x}(C)\Longrightarrow
  \operatorname{ExactlyOne}\bigl(\mbox{\CRETE}_{x}(C),\mbox{\CREBF}_{x}(C),\mbox{\CRECTMT}_{x}(C)\bigr).
\]
It is not proved here.  Since the three states are defined by data
recorded in the carrier, (Cls) is a condition on carrier records, an
admissibility requirement on how attempts are recorded, rather than a conjecture about proofs.  It requires that every \Xequiv\ carrier
record exactly one failure, so a successful \Xequiv\ attempt cannot be
recorded at all; the foreclosure of this part is a consequence of this
recording convention, not evidence about the problems.  Since the
states can overlap (see above), they form, on the carriers that satisfy
the Coverage schema below, a cover rather than a partition.

\begin{theorem}[Carrier Dichotomy, conditional form]\label{thm:carrier-dichotomy}
Let \(\mathcal B\) be a sound set of accepted claims, \(x\) a target
claim, and \(C\) a carrier that satisfies \textup{(Cls)}.  Then:
\begin{enumerate}
\item if \(C\) is \nonXequiv, there is a claim \(c'\neq x\) with
  \((R_{C}=0)\leftrightarrow\llbracket c'\rrbracket\) and
  \(\operatorname{NativeOutcome}_{\mathcal B}(C,c')\): either \(c'\) is
  accepted and \(R_{C}=0\), or \(c'\) is open-pending;
\item if \(C\) is \Xequiv, exactly one of \CRETE, \CREBF\ and
  \CRECTMT\ holds.
\end{enumerate}
\end{theorem}

\begin{proof}
For (1), choose \(c'\) as in the definition.  If \(c'\in\mathcal B\),
soundness gives \(\llbracket c'\rrbracket\), hence \(R_{C}=0\); otherwise
\(c'\) is open-pending.  Part (2) is (Cls).
\end{proof}

The first branch is the only one proved from the definitions; the
second is the classification hypothesis.  The theorem needs no aiming,
typing or coverage-context hypothesis; these enter only through the
intended scope of (Cls).  The \CRETE\ alternative is
the case in which the attempt's same-level audit has failed in the
sense of \Cref{sec:apparatus:birdint}; \CREBF\ and \CRECTMT\ are the
bridge and cascade alternatives.

\subsection{Bridge Impossibility}\label{sec:foreclosure:bridge-impossibility}

A \nonXequiv\ carrier might be enriched by a bridge into a composite
that is \Xequiv.  The question is what status such a bridge has.

\begin{definition}[Carrier-side typed bridge]\label{def:carrier-side-typed-bridge}
A \emph{\TypedBridgeCarrier} is a directed transport from a carrier
toward the target side.  The primitive relation
\(\operatorname{BridgeComp}_{x}(B,C,C_{B})\) states that composing the
bridge \(B\) with the carrier \(C\) yields the composite carrier
\(C_{B}\) for the target \(x\).\footnote{This is distinct from the
calibration-side bridge of \Cref{def:calibration-side-typed-bridge}.}
\end{definition}

\begin{definition}[\(X\)-strength obligation]\label{def:x-strength-obligation}
The primitive predicate \(\operatorname{XObl}(B,x)\) states that the
bridge \(B\) carries an \emph{\XstrengthObl}: it must itself supply the
target-strength content needed for the composite to reach \(x\), and is
not merely a change of presentation.
\end{definition}

The \emph{bridge-obligation hypothesis} for \((B,C,C_{B},x)\) is
\[
\begin{aligned}
  \mbox{(BO)}\qquad
    &\operatorname{BridgeComp}_{x}(B,C,C_{B})\\
  &\quad\wedge\operatorname{ExactlyOne}\bigl(\mbox{\CRETE}_{x}(C_{B}),\mbox{\CREBF}_{x}(C_{B}),\mbox{\CRECTMT}_{x}(C_{B})\bigr)\\
  &\quad\Longrightarrow\operatorname{XObl}(B,x).
\end{aligned}
\]

\begin{corollary}[Bridge Impossibility, conditional form]\label{cor:bridge-impossibility}
Let \(\mathcal B\) be a sound set of accepted claims and \(x\) a target
claim.  Let \(C\) be a \nonXequiv\ carrier with \(R_{C}=0\), let \(B\) be
a carrier-side bridge, and let \(C_{B}\) be a composite with
\(\operatorname{BridgeComp}_{x}(B,C,C_{B})\), \(R_{C_{B}}=0\),
\((R_{C_{B}}=0)\leftrightarrow\llbracket x\rrbracket\),
\(C_{B}\leadsto x\), \(C_{B}\models\text{\TCD}\) and
\(\operatorname{CoverageCtx}(x,C_{B})\).  If \(C_{B}\) satisfies
\textup{(Cls)} and \textup{(BO)} holds, then \(\operatorname{XObl}(B,x)\).
\end{corollary}

\begin{proof}
The composite is \Xequiv, so \Cref{thm:carrier-dichotomy} applied to
\(C_{B}\) gives exactly one of its three carrier-residual-equivalent
states, and (BO) gives \(\operatorname{XObl}(B,x)\).
\end{proof}

The hypotheses of the corollary imply that \(\llbracket x\rrbracket\)
holds, since \(C_{B}\) is \Xequiv\ and \(R_{C_{B}}=0\); the corollary
describes what a bridge that actually reached the target would have to
carry.  The hypotheses are jointly satisfiable: there is a model in
which all of them hold at once, with two different claims that are both
true (\Cref{app:formalization}).  The model is degenerate (a
single carrier with zero matrices serves as both \(C\) and \(C_{B}\)),
so it shows only that the hypotheses do not contradict each other, not
that (BO) is plausible.  The proof uses only the
composition, the \(X\)-equivalence of \(C_{B}\), (Cls) and (BO); the
remaining hypotheses describe the situation the corollary is about.  The content of the corollary is
carried by (BO): a bridge that turns a \nonXequiv\ carrier into an
\Xequiv\ composite must carry the target-strength content it appears to
transport.  This is a statement inside the carrier setup of this part,
not a general impossibility result for bridges into \(x\).

\subsection{Coverage}\label{sec:foreclosure:coverage}

\begin{schema}[Coverage Theorem]\label{thm:coverage}
Let \(x\) be a target claim and \(C\) a carrier with \(C\leadsto x\),
\(C\models\text{\TCD}\) and \(\operatorname{CoverageCtx}(x,C)\).  Then
\[
  \mbox{\nonXequiv}_{x}(C)\vee
  \mbox{\CRETE}_{x}(C)\vee
  \mbox{\CREBF}_{x}(C)\vee
  \mbox{\CRECTMT}_{x}(C).
\]
\end{schema}

The schema asserts exhaustiveness only.  For a carrier that records a
claim \(r_{C}\) with \(\llbracket r_{C}\rrbracket\leftrightarrow(R_{C}=0)\),
the schema follows from (Cls): if \(C\) is not \Xequiv, then
\(r_{C}\neq x\) and \(r_{C}\) witnesses that \(C\) is \nonXequiv;
otherwise (Cls) gives one of the three states.  Without such a claim it
need not follow: in a language with a single claim, interpreted as
false, a carrier with \(R_{C}=0\) and no recorded failures satisfies
(Cls) vacuously but none of the four states.  Turning the schema into
a statement about actual proof attempts would require definitions of
the recorded predicates and a notion of a target that is open relative
to \(\mathcal B\), neither settled by accepted claims nor replaceable by
an equivalent claim that imports accepted content; this paper supplies
neither.

\subsection{Type-Finite Foreclosure}\label{sec:foreclosure:type-finite}

The set of possible carriers aimed at \(x\) may be infinite; the set of
audit outcomes is not.  Let
\(\mathsf{State}=\{\mathsf{nonX},\mathsf{TE},\mathsf{BF},\mathsf{CTMT}\}\),
and let \(\operatorname{ClassifiesAs}(s,x,C)\) mean that \(C\) is in the
state named by \(s\).

\begin{theorem}[Type-Finite Foreclosure, conditional form]\label{thm:type-finite-foreclosure}
If a carrier \(C\) satisfies the conclusion of
the Coverage schema for \(x\), then
\(\operatorname{ClassifiesAs}(s,x,C)\) for some \(s\in\mathsf{State}\).
Consequently, any family of carriers satisfying the Coverage schema for
\(x\) is classified by a map into the four-element set
\(\mathsf{State}\), however large the family is.
\end{theorem}

\begin{proof}
Choose the state named by a true disjunct of the coverage disjunction.
\end{proof}

Since the states can overlap, the classifying map is a choice and is
not canonical.

\subsection{Attack Foreclosure}\label{sec:foreclosure:attack-foreclosure-v4}

The framework has three kinds of \emph{\fim}: a \emph{\carrierPivot}
replaces the carrier by another carrier aimed at \(x\); a
\emph{\bridgeImport} composes the carrier with a bridge to \(x\); and a
\emph{\cascadeInternal} continues the reduction inside the carrier.
None of these moves is forbidden.  The claim of this subsection is that
they do not, by themselves, produce progress on \(x\) without content
from outside the framework.

\begin{definition}[Framework-scope predicate]\label{def:in-scope}
\(\operatorname{InScope}(x)\) is a primitive predicate marking the
targets to which the framework's typing and closure-formation discipline
applies.  \(\operatorname{AdvancesWithoutExternal}(m,x)\) is a
primitive predicate stating that the internal move \(m\) makes progress
on \(x\) using only the framework's internal operations.  The scope
predicate is meant to exclude targets already settled by accepted
claims, for which the schema below would fail under any natural
reading; the formal statement does not enforce this.
\end{definition}

\begin{schema}[Attack Foreclosure]\label{thm:attack-foreclosure-v4}
For every target \(x\) with \(\operatorname{InScope}(x)\) and every
framework-internal move
\(m\in\{\mbox{\carrierPivot},\mbox{\bridgeImport},\mbox{\cascadeInternal}\}\),
\[
  \neg\,\operatorname{AdvancesWithoutExternal}(m,x).
\]
\end{schema}

The intended reading is that each move either returns to the
four-state classification or exposes outside content that must be
accepted through closure formation.  This reading is an interpretation,
not a consequence of the schema: the predicates are primitive, so the
schema neither constructs a transition between carriers nor exhibits
the outside content.  The intended derivation is a case analysis.  A
pivot gives a new carrier aimed at \(x\), which
\Cref{thm:carrier-dichotomy} and the Coverage schema return to the four
states.  A bridge import that reaches an \Xequiv\ composite falls under
\Cref{cor:bridge-impossibility}.  A cascade-internal computation either
closes the residual with accepted content or enters a tower of
matrix-element terminations (\Cref{thm:ctmt-recursion}).  For a finite
chain of moves the same alternative holds move by move.  What is
missing is the theorem that this case analysis covers every admissible
chain of internal moves, which is why the statement is a schema.

The schema concerns the framework's own moves.  It is not a barrier
theorem in the sense of the classical results: relativization, by
oracles relative to which \(P=NP\) and others relative to which
\(P\neq NP\) \cite{BakerGillSolovay1975Relativization}; natural proofs,
which cannot prove superpolynomial circuit lower bounds if suitable
pseudorandom functions exist \cite{RazborovRudich1997NaturalProofs};
and algebrization, an extension of relativization to algebraic oracles
\cite{AaronsonWigderson2009Algebrization}.  The conditional results on feature-based proof
attempts in the P-versus-NP application \cite{Tsiokos2026PvNP} are
likewise conditional on explicit hypotheses and are not a barrier of
that scope.

\subsection{CTMT Recursion}\label{sec:foreclosure:ctmt-recursion}

The cascade-terminal state need not be a single step: the decision of a
terminal matrix family may itself reduce to a deeper family.  Fix a
relation \(\operatorname{Reduces}_{\sigma}(x,C,C')\), the \emph{cascade
semantics}, stating that the cascade of \(C\) reduces to that of \(C'\).
Define towers inductively: \(\mathsf{Tower}_{\sigma}(x,C,1)\) holds when
\(C\) is in \CRECTMT, and \(\mathsf{Tower}_{\sigma}(x,C,n+2)\) holds when
\(\mathsf{Tower}_{\sigma}(x,C',n+1)\) and \(\operatorname{Reduces}_{\sigma}(x,C,C')\)
for some \(C'\).

\begin{theorem}[CTMT Recursion]\label{thm:ctmt-recursion}
For every cascade semantics \(\sigma\), target \(x\) and carrier \(C\),
\[
  \mathsf{Tower}_{\sigma}(x,C,1)\Longleftrightarrow\mbox{\CRECTMT}_{x}(C),
\]
and for every \(n\ge0\), \(\mathsf{Tower}_{\sigma}(x,C',n+1)\) and
\(\operatorname{Reduces}_{\sigma}(x,C,C')\) imply
\(\mathsf{Tower}_{\sigma}(x,C,n+2)\).
\end{theorem}

\begin{proof}
A tower of depth \(1\) can only come from the first clause, since the
second produces depth at least \(2\); conversely the first clause
applies to every carrier in \CRECTMT.  The second statement is the
second clause.
\end{proof}

The theorem records that matrix-element termination can occur at any
finite depth along a chain of reductions, without introducing a fifth
state.  It does not construct a reduction in any application; the
relation \(\operatorname{Reduces}_{\sigma}\) is supplied.

\begin{table}[H]
  \centering
  \caption{The statements of Part~I and what each rests on.}
  \label{tab:foreclosure:summary}
  \small
  \renewcommand{\arraystretch}{1.1}
  \begin{tabularx}{\textwidth}{@{}>{\raggedright\arraybackslash}p{0.25\textwidth}>{\raggedright\arraybackslash}X>{\raggedright\arraybackslash}p{0.27\textwidth}@{}}
    \toprule
    Statement & Content & Status \\
    \midrule
    \Cref{thm:carrier-dichotomy} & native-outcome branch for \nonXequiv\ carriers; exactly one failure state for \Xequiv\ ones & proved, assuming (Cls) \\
    \Cref{cor:bridge-impossibility} & a bridge into an \Xequiv\ composite carries an \XstrengthObl & proved, assuming (Cls) and (BO) \\
    \Cref{thm:coverage} & the four states are exhaustive & schema; follows from (Cls) when some claim expresses \(R_{C}=0\) \\
    \Cref{thm:type-finite-foreclosure} & audit outcomes form a four-element set & proved, assuming coverage \\
    \Cref{thm:attack-foreclosure-v4} & internal moves do not advance an in-scope target without outside content & schema \\
    \Cref{thm:ctmt-recursion} & matrix-element termination nests to any finite depth & proved \\
    \bottomrule
  \end{tabularx}
\end{table}


\section{Part II --- Recognition-mode landings}\label{sec:recognition}

\subsection{Two columns}\label{sec:recognition:column-classification}

Part~I described how internal derivations stall.  Part~II describes the
other route used by the companion applications: a closure that names
its load-bearing source as a premise, instead of presenting it as the
output of an internal derivation.  We call such a closure a
\emph{\RML}.  This part is a vocabulary together with four named
patterns, stated as conditional schemas.  Their predicates---whether a
template applies, whether content is derivable from the framework's
primitives, whether a closure lands at theorem grade---are primitive, so
the schemas have no truth conditions until those predicates are given a
meaning; they describe an intended informal reading and are not
mathematical conjectures in the strict sense.  The two applications
that illustrate them are described in \Cref{sec:validation}.  Unlike
Part~I, the targets here are propositions rather than claims; only
\Cref{thm:attack-foreclosure-v5}, which refers back to Part~I, uses the
interpretation of a claim.

The organizing distinction comes from the V-Differential classification
of the hiddenness companion \cite{Tsiokos2026CSL}, which asks how a
substrate's witness is exposed to current observation.  In a
witness-content-exposing substrate, current observation can read a
witness datum, such as an assignment or a solution coordinate.  In a
trace-state-only substrate, current observation sees only trace
bookkeeping---step counts, audit fields and similar records of how a
computation proceeded---and no witness datum.

\begin{definition}[Columns]\label{def:vdiff-tso-column}
A substrate \(S\) lies in the \emph{trace-state-only column},
\(\operatorname{TSO}(S)\), when its load-bearing witness is not exposed to
current observation and is available only through trace-state evidence,
which a closure must name and accept as a premise.  It lies in the
\emph{witness-content-exposing column}, \(\operatorname{WCE}(S)\), when its
witness is exposed to current observation through the substrate's own
apparatus.  Both are primitive predicates.
\end{definition}

The hiddenness companion places the Navier--Stokes substrate in the
second column and the Riemann-hypothesis and P-versus-NP substrates in
the first.  It adopts these placements as hypotheses, with an informal
motivation for each, and this paper uses them with the same status.
For the Riemann hypothesis the named source is
\(\Gamma_{\mathrm{SDTC}}^{\mathrm{Selberg}}\), recorded in the
application \cite{Tsiokos2026RH}; the self-dual trace confinement
companion \cite{Tsiokos2026SDTC} supplies the squeeze it feeds.  Its name refers to Selberg's trace
formula \cite{Selberg1956TraceFormula} by analogy only: Selberg's
formula concerns automorphic spectra and does not provide the
completed-zeta data that the source asserts, which remain a premise.
For \(P\) versus \(NP\) the named source is
\(\Gamma_{\mathrm{csl\text{-}sat\text{-}hidden}}\), recorded in the
application \cite{Tsiokos2026PvNP}.

Naming a source is not a sufficient premise in itself.  The
Riemann-hypothesis endpoint also needs an identification of its
coordinates with the classical zeros, including the conditions that
the identification preserves the real part \(\frac12\) and covers every
zero; the Navier--Stokes endpoint needs its spectral closure and its
hypotheses on the initial data; and the P-versus-NP endpoint needs an
accepted negative source and, in its family form, the admission of
search events.  Recognition-mode landings and derivation-mode landings
are two different kinds of conditional closure; neither is a weaker
form of the other.

\subsection{The seven-stage template}\label{sec:recognition:template-applicability}

The template describes a recognition-mode landing as a sequence of
seven checks on a formed closure \(c\) with target \(X\).
\begin{enumerate}
\item \emph{Construction.}  Build the formed closure with its types,
  levels, thresholds and visibility data, and exhibit that it reaches
  beyond the carriers available before it was formed.
\item \emph{Translation.}  Prove that the closure's internal predicate
  is equivalent to the standard statement of \(X\).  This is not where
  the source enters; it ensures that an accepted source yields the
  standard statement.
\item \emph{Applicability.}  Check that the relevant structural theorem
  applies to the closure, and separate the hypotheses already verified
  from those still needed.
\item \emph{Smuggling check.}  Check that the admissibility,
  lawfulness and saturation hypotheses do not already assume the source.
\item \emph{Anti-tautology.}  Check that the closure's predicate
  genuinely depends on the closure's data, so that the apparent gain
  does not split into independent claims about independent carriers.
\item \emph{Derivation sweep.}  Search for a derivation of the source
  along three routes (\Cref{def:three-option-sweep}).
\item \emph{Recognition.}  Name the source and record its acceptance as
  a separate premise.
\end{enumerate}
We write \(T=\langle S_{1},\ldots,S_{7}\rangle:\mathsf{7Stage}(c,X)\) for a
record of the seven stages as propositions about \((c,X)\).

\begin{definition}[Template applies]\label{def:applies}
For a formed closure \(c\), a target \(X\) and a seven-stage record
\(T\), \(\operatorname{Applies}(c,X,T)\) is a primitive predicate,
intended to state that the seven checks of \(T\) succeed for \((c,X)\);
it is not defined as the conjunction of the seven stages.
\end{definition}

\begin{schema}[Recognition-Mode Template Applicability]\label{thm:template-applicability}
Let \(c\) be a formed closure with target \(X\), and suppose that
\(\operatorname{TSO}(c)\), that the SB closure assumption holds for \(c\),
and that \(T:\mathsf{7Stage}(c,X)\) satisfies
\(\operatorname{Applies}(c,X,T)\).  Then \(c\) lands at theorem grade:
\(\operatorname{TheoremGrade}(c,X)\).
\end{schema}

All four inputs are needed; in particular the schema assumes, and does
not derive, that the template applies.  \(\operatorname{TheoremGrade}\)
is primitive, so even an instance of the schema would not prove \(X\);
it would record that the closure has the status of a conditional
theorem under its premises.  The two applications of
\Cref{sec:validation} illustrate the schema.

\subsection{The derivation sweep}\label{sec:recognition:verdict-signature}

\begin{definition}[Three-option derivation sweep]\label{def:three-option-sweep}
The \emph{derivation sweep} of Stage~6 tests three routes to the
load-bearing source: the \emph{framework-primitive option}, a
derivation from the primitives of the framework; the \emph{non-descent
option}, whether a derivation is blocked because the needed content
does not descend to the coarser layer; and the \emph{elevation option},
whether passing to the V-Differential classification resolves the
route.  Each option receives one of three verdicts: (i) the option is
the closure route; (ii) the option is not the closure route but does
not block it; (iii) the option passes the question to the
V-Differential classification.
\end{definition}

Formally, \(\mathsf{SweepOption}=\{\mathsf{FP},\mathsf{ND},\mathsf{VDE}\}\),
\(\mathsf{Verdict}=\{\mathrm{i},\mathrm{ii},\mathrm{iii}\}\), and a
seven-stage record \(T\) has a primitive verdict map
\(v_{T}:\mathsf{SweepOption}\to\mathsf{Verdict}\).

\begin{schema}[Cross-Track Verdict Signature]\label{thm:verdict-signature}
Let \(c\) be a formed closure in the trace-state-only column with
target \(X\), and let \(T:\mathsf{7Stage}(c,X)\) satisfy
\(\operatorname{Applies}(c,X,T)\).  Then
\[
  v_T(\mathsf{FP})=\mathrm{ii},\qquad
  v_T(\mathsf{ND})=\mathrm{iii},\qquad
  v_T(\mathsf{VDE})=\mathrm{ii}.
\]
\end{schema}

The intended explanation is that a verdict (i) for the
framework-primitive option would contradict
\Cref{thm:closure-content-thesis}, and that the other two options only
relabel the route; this is a heuristic, not a proof.

The evidence is limited to the two trace-state-only applications.  The
first version of the Riemann-hypothesis application reports that three
candidate derivations of its source from the framework's primitives
each fail \cite{Tsiokos2026RHvOne}, and the first version of the
P-versus-NP application reports the same for its source
\cite{Tsiokos2026PvNPvOne}.
\Cref{tab:verdict-signature} records this paper's reading of those
reports as verdicts.  The assignment is not computed by any formal
procedure, and no formal instance of the schema is constructed.

\begin{table}[H]
  \centering
  \caption{Verdicts of the derivation sweep, as read in this paper from
  the reports of the two applications; not a formal computation.}
  \label{tab:verdict-signature}
  \small
  \begin{tabularx}{\textwidth}{@{}>{\raggedright\arraybackslash}p{0.22\textwidth}>{\centering\arraybackslash}X>{\centering\arraybackslash}X>{\centering\arraybackslash}X@{}}
    \toprule
    Application & Framework-primitive option & Non-descent option & Elevation option \\
    \midrule
    \(P\) versus \(NP\) & (ii) & (iii) & (ii) \\
    Riemann hypothesis & (ii) & (iii) & (ii) \\
    \bottomrule
  \end{tabularx}
\end{table}


\subsection{The closure-content thesis}\label{sec:recognition:closure-content-thesis}

The derivation sweep suggests a semantic distinction.  On a
trace-state-only substrate the load-bearing content is part of the
formed closure itself: the closure names it and accepts it as a
premise.  We use three primitive predicates of a structural fact
\(s\): \(\operatorname{FLC}(s,c)\), that \(s\) is \emph{formed-layer
content} of \(c\), supplied by the formed closure rather than produced
later; \(\operatorname{DFP}(s)\), that \(s\) is \emph{derivable from
the framework's primitives}; and \(\operatorname{NAS}(s,c)\), that \(s\)
is \emph{named as an accepted source} of \(c\).

\begin{definition}[Load-bearing content]\label{def:load-bearing}
For a formed closure \(c\) and a target \(X\), the load-bearing content
\(s_{c,X}=\operatorname{LoadBearingFor}(c,X)\) is the structural fact
that \(c\) must accept as a premise in order to land \(X\).  The
assignment \((c,X)\mapsto s_{c,X}\) is primitive.
\end{definition}

\begin{schema}[Closure-Content Thesis]\label{thm:closure-content-thesis}
Let \(c\) be a formed closure in the trace-state-only column, with
target \(X\).  Then its load-bearing content is formed-layer content,
is not derivable from the framework's primitives, and is named as an
accepted source:
\[
  \operatorname{FLC}(s_{c,X},c)\wedge
  \neg\operatorname{DFP}(s_{c,X})\wedge
  \operatorname{NAS}(s_{c,X},c).
\]
\end{schema}

The three predicates are independent: a fact can be formed-layer
content without being named, be named without being load-bearing, or be
derivable without being a recognition source.  The thesis asserts the
specific combination for the load-bearing content of a trace-state-only
closure.  It concerns the content \(s_{c,X}\) only; that every
recognition-mode landing of \(c\) imports this content is a separate
clause of \Cref{thm:attack-foreclosure-v5}.

\subsection{Refined attack foreclosure}\label{sec:recognition:attack-foreclosure-v5}

The last schema combines the two parts.  Write
\(\operatorname{DerivationRoute}(m)\) for the primitive predicate that
an internal move \(m\) is a derivation route,
\(l:\operatorname{RMLanding}(c,X)\) for a recognition-mode landing of
\(X\) on \(c\), and \(\operatorname{ImportsAsNamedRecognition}(l,s)\)
for the primitive predicate that \(l\) imports \(s\) as a named source.

\begin{schema}[Refined Attack Foreclosure]\label{thm:attack-foreclosure-v5}
Let \(x\) be a claim in framework scope, put \(X=\llbracket x\rrbracket\),
and let \(c\) be a formed closure in the trace-state-only column for which
the SB closure assumption holds.  Then
\[
\begin{aligned}
  &\forall m.\;
    \operatorname{DerivationRoute}(m)
    \Longrightarrow
        \neg\operatorname{AdvancesWithoutExternal}(m,x),\\
  &\forall l:\operatorname{RMLanding}(c,X).\;
    \operatorname{ImportsAsNamedRecognition}(l,s_{c,X}),\\
  &\forall l:\operatorname{RMLanding}(c,X).\;
    \operatorname{TheoremGrade}(c,X).
\end{aligned}
\]
\end{schema}

The first clause is a consequence of \Cref{thm:attack-foreclosure-v4},
restricted to derivation routes.  The second and third clauses are new:
a recognition-mode landing imports the load-bearing content as a named
source, and lands at theorem grade.  The schema therefore adds to
\Cref{thm:attack-foreclosure-v4} and does not retract it.  It does not
assert that every admissible route is either a derivation route or a
recognition-mode landing; that the two kinds exhaust the routes is an
additional assumption, which we do not make.

\begin{table}[H]
  \centering
  \caption{The four statements of Part~II.  All are conditional schemas
  with primitive predicates; none is proved.}
  \label{tab:recognition:summary}
  \small
  \renewcommand{\arraystretch}{1.1}
  \begin{tabularx}{\textwidth}{@{}>{\raggedright\arraybackslash}p{0.27\textwidth}>{\raggedright\arraybackslash}X@{}}
    \toprule
    Statement & Content \\
    \midrule
    \Cref{thm:template-applicability} & On a trace-state-only formed closure satisfying the SB closure assumption, a seven-stage record for which the template applies yields theorem-grade landing. \\
    \Cref{thm:verdict-signature} & Under the same hypotheses, the derivation sweep returns the verdicts \SigPattern. \\
    \Cref{thm:closure-content-thesis} & The load-bearing content of a trace-state-only closure is formed-layer content, is not derivable from the primitives, and is named as an accepted source. \\
    \Cref{thm:attack-foreclosure-v5} & For an in-scope target and a trace-state-only formed closure satisfying the SB closure assumption: derivation routes do not advance the target without outside content, and recognition-mode landings import the load-bearing content and land at theorem grade. \\
    \bottomrule
  \end{tabularx}
\end{table}


\section{The three applications}\label{sec:validation}

\subsection{Overview}\label{sec:validation:intro}

This section places the three companion applications in the scheme of
\Cref{sec:recognition}.  It derives nothing: the applications are cited
as they stand, and no theorem of theirs is restated or re-proved.  The
problems themselves are stated in the Clay Mathematics Institute
descriptions \cite{Fefferman2000NavierStokes,Bombieri2000Riemann,Cook2000PvsNP};
for recent work on Navier--Stokes regularity see
\cite{Tao2016NSAveragedBlowup,BuckmasterVicol2019NS}, for the Riemann
hypothesis \cite{Conrey2003RH,IwaniecKowalski2004ANT}, and for
NP-completeness \cite{Cook1971NPCompleteness,Karp1972NPComplete21}.

The three applications do not share one closure hypothesis.  The
Navier--Stokes application \cite{Tsiokos2026NavierStokes} assumes a
spectral closure hypothesis, uniform over the time horizon, together with hypotheses on
real, divergence-free Sobolev initial data.  The Riemann-hypothesis
application \cite{Tsiokos2026RH} assumes a supplied Selberg recognition
source and an identification with the classical zeros of the completed
zeta function.  The P-versus-NP application \cite{Tsiokos2026PvNP}
proves a fixed machine-level translation and derives \(P\neq NP\) from
an accepted negative source, whose content and, in the family form,
whose admission to the chosen family of observables are premises.  The
SB closure assumption discharges none of these.

The applications also fall on both sides of the column distinction of
\Cref{sec:recognition:column-classification}.  As placed in the
hiddenness companion \cite{Tsiokos2026CSL}, the Navier--Stokes
substrate is witness-content-exposing, and the Riemann-hypothesis and
P-versus-NP substrates are trace-state-only, with the recognition
sources \(\Gamma_{\mathrm{SDTC}}^{\mathrm{Selberg}}\) and
\(\Gamma_{\mathrm{csl\text{-}sat\text{-}hidden}}\) respectively.  The
placements are hypotheses of that paper, and the sources are premises
of the applications; neither is derived here.

\subsection{Navier--Stokes}\label{sec:validation:ns}

The Navier--Stokes application is the witness-content-exposing case.
Its conditional closure proceeds through problem-specific analytic
apparatus, including the layer-dissolving membrane of
\cite{Tsiokos2026NeedleKiller,Tsiokos2026NavierStokes}, and not through
the recognition-mode template.  It therefore shows that the template
does not cover every closure in the series.  It is not evidence for the
schemas of \Cref{sec:recognition}; under the placements of the
hiddenness companion, it is a case on the other side of the column
distinction.

\subsection{The Riemann hypothesis}\label{sec:validation:rh}

The Riemann-hypothesis application is the first trace-state-only case.
Its load-bearing source is the Selberg recognition source
\(\Gamma_{\mathrm{SDTC}}^{\mathrm{Selberg}}\), recorded in the
application \cite{Tsiokos2026RH} and not derived from framework
primitives.  In the terms of \Cref{sec:recognition}, the application
follows the seven stages, the derivation sweep reported in its first
version \cite{Tsiokos2026RHvOne} returns, in this paper's reading of
that report, the verdicts \SigPattern, its load-bearing source is named rather than derived, and
its route is a recognition-mode landing rather than a derivation.  These
four observations are what the application contributes as an
illustration of the four schemas; they are read off the application,
not proved from the schemas.

\subsection{P versus NP}\label{sec:validation:pvnp}

The P-versus-NP application works with the classes \(P\) and \(FP\)
defined by deterministic multi-tape Turing machines with polynomial
running time, the class \(NP\) defined by nondeterministic ones, and
the standard bitstring encoding of
CNF formulas, as formalized in complexitylib
\cite{Schlesinger2026Complexitylib}.  Its \emph{tagged witness
selector} sends a bitstring \(z\) that encodes a satisfiable formula to
the tag bit \(1\) followed by the lexicographically first satisfying
assignment of the variables \(0,\ldots,|z|\), and every other bitstring,
including malformed ones, to the single bit \(0\).  The application proves, in both directions, that a
polynomial-time SAT decider exists exactly when this selector is
computable in polynomial time.  If the selector is not computable in polynomial time,
then \(P\neq NP\); this negative premise is not proved.  The family
form of the result additionally requires that every constructed search
event be admitted to the chosen family of current observables
\cite{Tsiokos2026PvNP}.

The proposed recognition source is
\(\Gamma_{\mathrm{csl\text{-}sat\text{-}hidden}}\), recorded in the
P-versus-NP paper \cite{Tsiokos2026PvNP}.  Its negative content and its acceptance are
separate premises; under the admission law, the accepted negative
content has the strength of \(L_{\mathrm{SAT}}\notin P\), which by the
Cook--Levin theorem \cite{Cook1971NPCompleteness} is equivalent to
\(P\neq NP\).  It is not
derived here, nor by the Gaussian construction of
\Cref{sec:triad-bridge}.  As for the Riemann hypothesis, the
application follows the seven stages, the sweep reported in its first
version \cite{Tsiokos2026PvNPvOne} returns, in this paper's reading of
that report, the verdicts \SigPattern, and its source is named rather than derived; it is the
second illustration of the schemas of \Cref{sec:recognition}.  The
conditional feature-decomposition results of the P-versus-NP paper
reappear in \Cref{sec:calibration:functorial-gates}.

\begin{table}[H]
  \centering
  \caption{The three applications.  Column placements are hypotheses of
  the hiddenness companion; sources and premises are those of the
  applications.}
  \label{tab:cross-track-validation}
  \small
  \renewcommand{\arraystretch}{1.1}
  \begin{tabularx}{\textwidth}{@{}>{\raggedright\arraybackslash}p{0.15\textwidth}>{\raggedright\arraybackslash}p{0.19\textwidth}>{\raggedright\arraybackslash}p{0.22\textwidth}>{\raggedright\arraybackslash}X@{}}
    \toprule
    Application & Column & Route & Premises \\
    \midrule
    Navier--Stokes & witness-content-exposing & problem-specific analytic apparatus & spectral closure; real divergence-free initial data of Sobolev index \(\gamma>5/2\); positive viscosity \\
    Riemann hypothesis & trace-state-only & recognition-mode, source \(\Gamma_{\mathrm{SDTC}}^{\mathrm{Selberg}}\) & accepted source; identification with the classical zeros that preserves the real part \(\frac12\) and covers every zero \\
    \(P\) versus \(NP\) & trace-state-only & fixed machine translation; family form recognition-mode, source \(\Gamma_{\mathrm{csl\text{-}sat\text{-}hidden}}\) & tagged witness selector not in \(FP\); in the family form, an accepted negative source and admission of search events \\
    \bottomrule
  \end{tabularx}
\end{table}


\section{A Gaussian observation shared by the Riemann and Navier--Stokes settings}
\label{sec:gaussian-bridge}

The comparisons of the preceding sections are structural.  This section
and the next record a concrete construction instead: one family of
complex Gaussian observations is applied both to finite sets of zeros of
the completed zeta function and to the Fourier side of the
three-dimensional Navier--Stokes equations, and in
\Cref{sec:triad-bridge} also to the assignments of a CNF formula.  The
identities are elementary: they come from completing the square in a
Gaussian integral and from the linearity of the observation.  They are
proved without any of the conditional schemas of
\Cref{sec:foreclosure,sec:recognition}, and their formal proofs use
only the standard axioms (\Cref{app:formalization}).  What they show is
that a single observation family can be read on all three objects; they
do not show that the hypotheses of the three conditional closures
support one another.

\subsection{The observation}

Write \(\xi\in\mathbb R^{3}\) for the Fourier variable.  Following the
Navier--Stokes companion \cite{Tsiokos2026NavierStokes}, fix a Sobolev
index \(\gamma>5/2\) and represent a velocity field \(u\) by its
\emph{weighted Fourier state}
\(g(\xi)=(1+|\xi|^{2})^{\gamma/2}\hat u(\xi)\), an element of
\(L^{2}(\mathbb R^{3};\mathbb C^{3})\); the physical amplitude is
\(\hat u=(1+|\xi|^{2})^{-\gamma/2}g\).  The observations below are
applied to weighted states.  For a width
\(a>0\), a unit vector \(\theta\in\mathbb R^{3}\) and a complex centre
\(z\in\mathbb C\), put
\[
  w_{a,\theta,z}(\xi)
  =\exp\bigl(-a|\xi|^{2}+2az\langle\theta,\xi\rangle-az^{2}\bigr),
  \qquad
  \mathcal O_{a,\theta,z}(g)=\int_{\mathbb R^{3}}w_{a,\theta,z}(\xi)\,g(\xi)\,d\xi .
\]
For real \(z\) the window is the Gaussian
\(\exp(-a|\xi-z\theta|^{2})\) centred at \(z\theta\); complex \(z\) is
its analytic continuation in the centre.  We fix \(\theta=\theta_{\rm
axis}\), the first coordinate vector, whenever a single direction is
needed.

Let \(\nu>0\) be the viscosity and \(t>0\) a time, and put
\(b=4\pi^{2}\nu t\).  With the convention
\(\hat g(\xi)=\int g(x)e^{-2\pi i x\cdot\xi}\,dx\), the heat semigroup
\(e^{\nu t\Delta}\) acts on Fourier states, weighted or not, as
multiplication by \(e^{-b|\xi|^{2}}\).  Completing the square gives,
with \(\kappa=ab/(a+b)\),
\begin{equation}
\label{eq:joint-heat-kernel}
  N_{a,b}(z):=\int_{\mathbb R^{3}}w_{a,\theta,z}(\xi)\,e^{-b|\xi|^{2}}\,d\xi
  =\Bigl(\frac{\pi}{a+b}\Bigr)^{3/2}e^{-\kappa z^{2}},
\end{equation}
and, for every Fourier state \(g\),
\[
  \mathcal O_{a,\theta,z}\bigl(e^{\nu t\Delta}g\bigr)
  =e^{-\kappa z^{2}}\,
    \mathcal O_{a+b,\,\theta,\,az/(a+b)}(g).
\]

\subsection{The zero side}

Let \(\rho\) be a zero of the completed zeta function in the critical
strip \(0<\operatorname{Re}\rho<1\).  Its \emph{spectral coordinate} is
\[
  z_{\rho}=\operatorname{Im}\rho+i\bigl(\tfrac12-\operatorname{Re}\rho\bigr),
  \qquad\text{so that}\qquad
  \operatorname{Re}\bigl(z_{\rho}^{2}\bigr)
  =(\operatorname{Im}\rho)^{2}-\bigl(\operatorname{Re}\rho-\tfrac12\bigr)^{2}.
\]
By \eqref{eq:joint-heat-kernel}, the logarithmic excess
\[
  \ell_{a,b}(\rho)
  :=\log\frac{|N_{a,b}(z_{\rho})|}{(\pi/(a+b))^{3/2}}
        +\kappa\,(\operatorname{Im}\rho)^{2}
  =\kappa\,\bigl(\operatorname{Re}\rho-\tfrac12\bigr)^{2}
\]
is defined so as to isolate the horizontal distance of \(\rho\) from
the critical line.

The map \(J(\rho)=1-\bar\rho\) is the reflection in the critical line;
it sends zeros to zeros by the functional equation and the reality of
the completed zeta function on the real axis.  A \emph{reflected
window} is a finite set \(W\) of distinct zeros in the strip with
\(J(W)=W\); multiplicities enter only through the weights.
Let \(m:W\to\mathbb N\) be weights with \(m\circ J=m\), and put
\[
  E^{\mathrm{heat}}_{a,\nu,t}(W)=\sum_{\rho\in W}m(\rho)\,\ell_{a,b}(\rho),
  \qquad
  v_{W}(\rho)=\sqrt{m(\rho)}\,\bigl(\operatorname{Re}\rho-\tfrac12\bigr),
  \qquad
  L_{W}x=\sum_{\rho\in W}x_{\rho}
\]
for \(x\in\mathbb R^{W}\).  In the language of the adequacy-residual
companion \cite{Tsiokos2026XiCompanion}, the residual of the identity
target against the probe \(L_{W}\), at identity audit energy, is the
orthogonal projection \(\Xi(L_{W})\) of \(\mathbb R^{W}\) onto
\(\ker L_{W}\).  The reflection reverses the sign of each displacement
and preserves the weights, so \(L_{W}v_{W}=0\) and
\(\Xi(L_{W})v_{W}=v_{W}\).  Hence
\begin{equation}
\label{eq:joint-heat-xi}
  E^{\mathrm{heat}}_{a,\nu,t}(W)
    =\kappa\sum_{\rho\in W}m(\rho)\bigl(\operatorname{Re}\rho-\tfrac12\bigr)^{2}
  =\frac{ab}{a+b}\,\bigl\|\Xi(L_{W})v_{W}\bigr\|^{2}.
\end{equation}
We call \(\|\Xi(L_{W})v_{W}\|^{2}\) the \emph{residual charge} of the
window; it vanishes exactly when every zero of positive weight in \(W\)
lies on the critical line.  The first equality holds for any finite set
of points of the strip with any nonnegative weights.  On the zero side
the residual projection does no work: the reflection hypotheses place
\(v_{W}\) in \(\ker L_{W}\), so the residual charge is just
\(\|v_{W}\|^{2}\).  The projection matters on the SAT side of
\Cref{sec:triad-bridge}, where native and residual channels differ.

\subsection{The fluid side}

Let \(g_{0}\) be a weighted Fourier state.  A time
\(t\ge0\) is \emph{admissible} for \(g_{0}\) if the mild formulation of
the Navier--Stokes equations with data \(g_{0}\) has a solution on
\([0,t]\) in the class of weighted states of index \(\gamma\) used in
\cite{Tsiokos2026NavierStokes}; when the physical amplitude of \(g_{0}\)
is that of a real, divergence-free field, positive admissible times
exist by local existence.  On admissible times the maximal mild solution
satisfies, in weighted coordinates,
\[
  g(t)=e^{\nu t\Delta}g_{0}-\mathcal D[g](t),
\]
where \(\mathcal D[g](t)\) is the nonlinear heat--Duhamel term: the
integral over \([0,t]\) of the heat semigroup applied to the
Leray-projected convection term \(P(u\cdot\nabla)u\), written in the
weighted coordinates and with the normalization of the Navier--Stokes
companion \cite{Tsiokos2026NavierStokes}, where the convection term is
carried at Sobolev weight \(\gamma-1\) and the heat semigroup supplies
the missing derivative.  Applying the linear observation to this equation
gives
\begin{equation}
\label{eq:joint-navier-return}
  \mathcal O_{a,\theta,z}(g(t))
  +\mathcal O_{a,\theta,z}(\mathcal D[g](t))
  =\mathcal O_{a,\theta,z}(e^{\nu t\Delta}g_{0}).
\end{equation}
This is the mild equation, observed: it keeps the nonlinear term and
says nothing about its size.

\begin{theorem}[Joint Gaussian observation and source return]
\label{thm:joint-gaussian-return}
Let \(a>0\), \(\nu>0\), \(\gamma>5/2\), let \(g_{0}\) be a weighted
Fourier state, and let \(t>0\) be admissible for \(g_{0}\); put
\(b=4\pi^{2}\nu t\) and \(\kappa=ab/(a+b)\).  Let \(W\) be a reflected
window with weights \(m:W\to\mathbb N\) satisfying \(m\circ J=m\), and
let \(\rho\in W\).  Then:
\begin{enumerate}
\item the heat excess of the window is the residual charge times
  \(\kappa\), as in \eqref{eq:joint-heat-xi};
\item the heat observation at the zero has the value
  \(N_{a,b}(z_{\rho})=(\pi/(a+b))^{3/2}e^{-\kappa z_{\rho}^{2}}\);
\item the maximal mild solution satisfies \eqref{eq:joint-navier-return}
  with \(\theta=\theta_{\rm axis}\) and \(z=z_{\rho}\);
\item \(m(\rho)(\operatorname{Re}\rho-\frac12)^{2}\le
  \|\Xi(L_{W})v_{W}\|^{2}\).
\end{enumerate}
\end{theorem}

\begin{proof}
Parts (1) and (2) are \eqref{eq:joint-heat-kernel} and the computation
leading to \eqref{eq:joint-heat-xi}.  Part (3) is the observation of the
mild equation.  Part (4) holds because the charge is a sum of
nonnegative terms, one of which is \(m(\rho)(\operatorname{Re}\rho-\frac12)^{2}\).
\end{proof}

Only part (3) involves \(\gamma\), \(g_{0}\) and the admissibility of
\(t\); parts (1), (2) and (4) concern the scalar kernel and hold for
all \(a,\nu,t>0\).  The theorem collects four elementary facts; its
point is that one observation family is read on both objects.

\begin{figure}[tbp]
\centering
\begin{tikzpicture}[every node/.style={draw,rounded corners,
  align=center,inner sep=6pt,font=\small}]
\node[text width=7.0cm] (obs) at (0,2.0)
  {One family of complex Gaussian observations\\
   heat transform and residual projection};
\node[text width=6.5cm] (rh) at (-3.75,0)
  {Reflected windows of zeta zeros\\
   heat excess $=\frac{ab}{a+b}\times$ residual charge\\
   vanishing of the charge not supplied};
\node[text width=6.5cm] (ns) at (3.75,0)
  {Maximal mild Navier--Stokes solution\\
   observed equation keeps the Duhamel term\\
   global closure hypotheses not supplied};
\draw[->,thick] (obs) -- (rh);
\draw[->,thick] (obs) -- (ns);
\end{tikzpicture}
\caption{One observation construction read on two objects.  The lower
line of each box names what the identities do not provide.}
\label{fig:joint-observation-branches}
\end{figure}

The same observation can be applied after the high-frequency residual
projection of the Navier--Stokes companion, which gives a projected
form of \eqref{eq:joint-navier-return} with the projected Duhamel term
still present.  It can also be applied to three mild solutions at once,
one for each coordinate direction; this three-solution \emph{frame} is
the version used in \Cref{sec:triad-bridge}.

\subsection{What each conditional closure still needs}

For fixed \(a,\nu,t>0\), let \(\mathcal A_{a,\nu,t}\) be the statement
that the heat excess vanishes, with unit weights, on every canonical
window, the set of zeros in the strip at distance at most \(n\) from
\(\frac12\) (\(n\in\mathbb N\)).  These windows are reflected and
exhaust the zeros in the strip, so by \eqref{eq:joint-heat-xi}
\[
  \mathcal A_{a,\nu,t}
  \quad\Longleftrightarrow\quad
  \text{the Riemann hypothesis}.
\]
The equivalence is a reformulation, and it measures the strength of the
premise \(\mathcal A_{a,\nu,t}\): the premise is the Riemann hypothesis.
The Gaussian identities supply no prime-side or trace-formula argument
for it.  The Riemann-hypothesis companion \cite{Tsiokos2026RH} follows a
different conditional route, in which a supplied Selberg recognition
source together with an identification with the classical zeros implies
the Riemann hypothesis.

On the fluid side, the Navier--Stokes companion
\cite{Tsiokos2026NavierStokes} proves two conditional returns, for
positive viscosity and real divergence-free initial data in
\(H^\gamma(\mathbb R^3)\), \(\gamma>5/2\).  A spectral closure hypothesis
on the maximal mild solution gives a global lifespan and a classical
solution.  A physical-window closure, for some \(0<\varepsilon<1\), asks for
finite Fourier shell witnesses whose amplitude and reconstruction
bounds satisfy \(2^{-\varepsilon j}(R_{t,j}+e_{t,j})\le b_T(t)\) on the
shells of the Bradshaw--Grujić window, with a majorant
\(b_T\in L^{2/(1-\varepsilon)}(0,T)\) for every possible preterminal
horizon \(T\); it gives a global strong mild solution only
together with a terminal-divergence input of Bradshaw--Grujić type for
the same maximal path and windows; that input is assumed there, not
derived.  These hypotheses concern
high frequencies, reconstruction and long times; the zero-side premise
concerns horizontal displacement of zeros.  Neither is derived from the
other through the shared observation.  A common source theorem would
have to supply both, and none is proved here.

\section{A finite Gaussian and residual construction across all three applications}
\label{sec:triad-bridge}

The construction of \Cref{sec:gaussian-bridge} extends to a third
object, the assignment space of a CNF formula.  Let \(\varphi\) be a CNF
formula in \(n\) Boolean variables, let \(A_{n}=\{0,1\}^{n}\), and let
\(s_{\varphi}:A_{n}\to\{0,1\}\) be its assignment verifier,
\(s_{\varphi}(w)=1\) exactly when \(w\) satisfies \(\varphi\).  The
satisfiability language \(L_{\mathrm{SAT}}\) is taken in the bitstring
encoding of CNF formulas used for the machine-defined classes \(P\),
\(NP\) and \(FP\) in the P-versus-NP companion \cite{Tsiokos2026PvNP};
the verifier agrees with that encoding, in the sense that the encoding
of \(\varphi\) lies in \(L_{\mathrm{SAT}}\) exactly when some \(w\)
satisfies \(\varphi\).  We regard \(s_{\varphi}\) as a vector in the
\(2^{n}\)-dimensional Euclidean space \(\mathbb R^{A_{n}}\).

\subsection{The shared kernel and the SAT residual}

The finite form of the construction sums the heat observation
\(N_{a,b}\) of \eqref{eq:joint-heat-kernel} over a finite weighted set
of points.  On a reflected zero window the points are the spectral
coordinates of the zeros, with their weights.  On \(A_{n}\) the weight
of \(w\) is \(s_{\varphi}(w)\) and the point is
\(\frac12+i\operatorname{code}(w)\), where \(\operatorname{code}(w)\)
is the natural number with binary digits \(w\); its spectral coordinate
is the real number \(\operatorname{code}(w)\).  This choice of points is
a convention.  For real \(\kappa\) the resulting mass is
\begin{equation}
\label{eq:triad-sat-mass}
 M_\kappa(\varphi)=\sum_{w\in A_n}s_\varphi(w)\,
       e^{-\kappa\operatorname{code}(w)^2}.
\end{equation}
Every term is nonnegative and the terms with \(s_{\varphi}(w)=1\) are
positive, so \(M_{\kappa}(\varphi)>0\) exactly when \(\varphi\) is
satisfiable.  The same holds for any \(\{0,1\}\)-valued filter in place
of \(s_{\varphi}\); what is specific to SAT is the verifier, its
agreement with the encoded language, and the clause and branching laws
below.

Adding a clause \(\chi\) removes exactly the assignments that satisfy
\(\varphi\) and fail \(\chi\):
\begin{equation}
\label{eq:triad-clause-residual}
 M_\kappa(\varphi\wedge\chi)
   =M_\kappa(\varphi)-
     \sum_{\substack{w\models\varphi\\w\not\models\chi}}
       e^{-\kappa\operatorname{code}(w)^2}.
\end{equation}

On the Navier--Stokes side we use a \emph{frame} of three mild
solutions.  Fix \(\gamma>5/2\) and \(\nu,t_{0}>0\), and for
\(i=1,2,3\) let \(g^{(i)}\) be the maximal mild solution, in the
weighted coordinates of \Cref{sec:gaussian-bridge}, with initial state
\(g^{(i)}_{0}(\xi)=e^{-4\pi^{2}\nu t_{0}|\xi|^{2}}P(\xi)e_{i}\), where
\(P(\xi)\) is the Leray projection symbol and \(e_{i}\) the \(i\)-th
coordinate vector; the physical initial amplitude is
\((1+|\xi|^{2})^{-\gamma/2}g^{(i)}_{0}(\xi)\), the transform of a real,
divergence-free field.  A time \(\tau\ge0\) is admissible for the frame
if it is admissible for all three data; a positive such time exists.
The \emph{frame trace} of the observation at \(\tau\) is
\(T_{a,z}(\tau)=\sum_{i}\bigl(\mathcal O_{a,\theta_{\rm axis},z}(g^{(i)}(\tau))\bigr)_{i}\),
and \(T^{\mathcal D}_{a,z}(\tau)\) is the same expression with the
Duhamel terms \(\mathcal D[g^{(i)}](\tau)\) in place of \(g^{(i)}(\tau)\).
Since the Leray symbol has trace \(2\), the mild equation gives
\[
  \tfrac12\bigl(T_{a,z}(\tau)+T^{\mathcal D}_{a,z}(\tau)\bigr)=N_{a,b}(z),
  \qquad b=4\pi^{2}\nu(\tau+t_{0}).
\]
We call the left-hand side the \emph{normalized frame value}.  Replacing
\(N_{a,b}\) by it in the definition of the logarithmic excess of
\Cref{sec:gaussian-bridge} gives the \emph{frame heat excess} of a
window, and, by the displayed identity, every heat-side identity
transfers to the frame with the Duhamel terms kept.  In particular
\eqref{eq:triad-clause-residual} holds with each term
\(e^{-\kappa\operatorname{code}(w)^{2}}\) replaced by the raw frame value
\(T_{a,z}+T^{\mathcal D}_{a,z}\) at \(z=\operatorname{code}(w)\), which is
\(2(\pi/(a+b))^{3/2}\) times that term.  These are exact identities,
not estimates of the nonlinear terms.

The common probe on a finite index set \(I\) is the sum of coordinates,
\(L_{I}(v)=\sum_{i\in I}v_{i}\), and the residual \(\Xi(L_{I})\) is the
orthogonal projection onto \(\ker L_{I}\); its complement
\(\Pi(L_{I})=1-\Xi(L_{I})\) is the projection onto the constant vectors,
the \emph{native} channel.  On a reflected zero window the probe of the
displacement vanishes, \(L_{W}(v_{W})=0\), so the displacement lies
entirely in the residual.  On the assignment space,
\(L_{A_{n}}(s_{\varphi})\) is the number of satisfying assignments,
which is positive exactly when the encoding of \(\varphi\) lies in
\(L_{\mathrm{SAT}}\).  Thus satisfiability is read in the native
channel.  The residual alone does not detect it: the residual of
\(s_{\varphi}\) vanishes both for the empty conjunction (every
assignment satisfies it) and for a formula containing the empty clause
(none does).

A unit clause \(x_{i}=\beta\) acts on states by the mask
\((M_{i,\beta}v)(w)=v(w)\) if \(w_{i}=\beta\) and \(0\) otherwise, so
that \(s_{\varphi\wedge(x_{i}=\beta)}=M_{i,\beta}s_{\varphi}\).  Since
\(\Xi=1-\Pi\), the commutator of the mask with the residual is
\[
  \Xi M_{i,\beta}s_{\varphi}-M_{i,\beta}\Xi s_{\varphi}
  =M_{i,\beta}\Pi s_{\varphi}-\Pi M_{i,\beta}s_{\varphi}.
\]
A one-variable example shows that the mask need not commute with
\(\Xi\).  No operation on zero windows is identified with the mask.

\begin{theorem}[Concrete Gaussian--residual triad]
\label{thm:concrete-triad-gaussian-xi}
Let \(\gamma>5/2\), \(\nu,t_{0}>0\) and \(a>0\), let \(\tau\ge0\) be
admissible for the frame, and put \(t=\tau+t_{0}\),
\(b=4\pi^{2}\nu t\) and \(\kappa=ab/(a+b)\).  Let \(W\) be a reflected
window with unit weights and \(\rho\in W\), and let \(\varphi\) be a CNF
formula in \(n\) variables, \(\chi\) a clause, \(i\le n\) and
\(\beta\in\{0,1\}\).  Then:
\begin{enumerate}
\item the frame heat excess of \(W\) equals
  \(\kappa\,\|\Xi(L_{W})v_{W}\|^{2}\), and \(L_{W}(v_{W})=0\);
\item the raw frame value \(T_{a,z}+T^{\mathcal D}_{a,z}\) summed over the
  satisfying assignments of \(\varphi\), at \(z=\operatorname{code}(w)\),
  is nonzero exactly when
  \(\|\Pi s_{\varphi}\|^{2}+\|\Xi s_{\varphi}\|^{2}>0\), that is, exactly
  when \(\varphi\) is satisfiable;
\item \(L_{A_{n}}(s_{\varphi})>0\) exactly when the encoding of
  \(\varphi\) lies in \(L_{\mathrm{SAT}}\);
\item adding \(\chi\) subtracts its failure term from the summed raw
  frame value, as in \eqref{eq:triad-clause-residual};
\item \(s_{\varphi\wedge(x_{i}=\beta)}=M_{i,\beta}s_{\varphi}\), and the
  commutator identity above holds.
\end{enumerate}
\end{theorem}

\begin{proof}
Part (1) is \Cref{thm:joint-gaussian-return} transferred to the frame by
the normalized frame identity, together with \(L_{W}v_{W}=0\).  For
part (2), the frame identity turns the summed raw value into
\(2(\pi/(a+b))^{3/2}M_{\kappa}(\varphi)\), which is nonzero exactly when
\(\varphi\) is satisfiable; by orthogonality,
\(\|\Pi s_{\varphi}\|^{2}+\|\Xi s_{\varphi}\|^{2}=\|s_{\varphi}\|^{2}\),
the number of satisfying assignments.  Part (3) is the agreement of the
verifier with the encoded language.  Parts (4) and (5) follow by
evaluating formulas assignment by assignment and from \(\Xi=1-\Pi\).
\end{proof}

Only part (1) uses the zero window, and the Navier--Stokes hypotheses
enter only through the frame identity, which contributes the constant
factor \(2(\pi/(a+b))^{3/2}\) to the SAT statements; parts (3) and (5)
are statements about finite vectors, and part (2) says that
\(s_{\varphi}\neq0\).

\begin{figure}[tbp]
\centering
\begin{tikzpicture}[every node/.style={draw,rounded corners,
  align=center,inner sep=5pt,font=\small}]
\node[text width=10.7cm] (top) at (0,2.2)
  {Finite complex Gaussian observation on one Navier--Stokes frame\\
   heat identity with the nonlinear Duhamel terms kept};
\node[text width=5.2cm] (rh) at (-3.25,0)
  {Reflected window of zeta zeros\\
   probe $=0$; residual charge is horizontal displacement\\
   vanishing of the charge not supplied};
\node[text width=5.2cm] (sat) at (3.25,0)
  {Assignment state of a CNF formula\\
   probe $>0\Leftrightarrow$ satisfiable; clause and branching laws\\
   running time and negative source not supplied};
\draw[->,thick] (top) -- (rh);
\draw[->,thick] (top) -- (sat);
\end{tikzpicture}
\caption{One observation and one Navier--Stokes frame read on two
finite states.  The probe readouts differ, and each application keeps
its own hypotheses.}
\label{fig:triad-observation-span}
\end{figure}

\subsection{The limits of the common construction}

The difference between the two probe readouts is a mathematical
constraint.  If \(\varphi\) is satisfiable, no map
\(F:\mathbb R^{W}\to\mathbb R^{A_{n}}\), linear or not, can satisfy
\(L_{A_{n}}(F(v))=L_{W}(v)\) for every \(v\) and \(F(v_{W})=s_{\varphi}\):
it would send probe value \(0\) to a positive probe value.  This
excludes only maps of that kind; it says nothing about other transports
between the two settings.

The positivity test \eqref{eq:triad-sat-mass}, evaluated as written,
inspects all \(2^{n}\) assignments, and no polynomial-time evaluation
or precision bound for the Gaussian mass is proved.  Its terms can be as
small as \(e^{-\kappa(2^{n}-1)^{2}}\), so deciding positivity from a
fixed-point approximation of the mass with absolute error below the
smallest term would need a number of bits exponential in \(n\).  The P-versus-NP companion proves separately that a
polynomial-time SAT decider exists exactly when its tagged witness
selector is computable in polynomial time; admitting the resulting
transducer to a chosen family of current observables is a further
premise, and under that admission the accepted negative source has the
strength of \(L_{\mathrm{SAT}}\notin P\).  The Gaussian and residual
identities supply none of this.

Likewise, the zero-window premise must make the residual charge vanish,
and the Navier--Stokes closure must control high frequencies,
reconstruction and long times.  \Cref{thm:concrete-triad-gaussian-xi}
connects three objects through one kernel and one residual calculus
with exact identities; it gives no implication from the hypothesis of
one application to that of another, and it is not a single sufficient
condition for all three problems.

\section{Part III --- The calibration family}\label{sec:calibration}

\subsection{Calibrations}\label{sec:calibration:intro}

A \emph{calibration} records a detector class, the kind of obstruction
it reads, together with the kind of carrier on which that detector is
meaningful.  A low-degree polynomial detector, a mutual-information
detector, a stable local algorithm and a Nullstellensatz certificate
have different native carriers and different notions of descent, and a
transport between two of them is a theorem, not a change of notation.
Part~III records these distinctions in a small typed vocabulary:
calibration records, a criterion for when a new record merely relabels
an old one, records of bridges between calibrations, a table of covered
regions, a four-axis comparison of types, and records of open
transports called \emph{functorial gates}.

The numbered results of this section are elementary facts about
finite records of labels: each one checks a property of declared finite
data, and the names they carry describe the intended use of the data
rather than the depth of the proof.  We state them because they make the catalogue precise and
because the functorial gates of \Cref{sec:calibration:functorial-gates}
are the open problems the catalogue points to.  They prove nothing
about how well a detector performs, about the existence of a bridge, or
about inclusions between complexity classes; region names such as
\(\mathsf{NC}^{1}\) are labels.

\subsection{The seed catalogue}\label{sec:calibration:structure}

\begin{definition}[Calibration]\label{def:calibration}
A \emph{calibration} is a record
\[
  \mathsf{Calibration}
  =
  \{\mathsf{name},\ \mathsf{detectorClass},\ \mathsf{nativeCarrierType}\}
\]
of three labels: its name, the class of detector it uses, and the type
of carrier on which the detector is meaningful.
\end{definition}

The seed catalogue \(\mathcal C_{0}\) is the list of thirteen records in
\Cref{tab:calibration:seed}; ``seed'' means that it is a starting list
to which further calibrations can be added.  It is introduced in this
paper; its
entries name detector classes that have been considered in connection
with the P-versus-NP application, and they are not results of that
application.  Background for the individual detector classes includes
\cite{Lovasz1978Kneser} for the topological line,
\cite{MulmuleySohoni2001GCT} for geometric complexity theory,
\cite{Krajicek2019ProofComplexity} for proof complexity and
Nullstellensatz certificates, and
\cite{AlonKrivelevichSudakov1998PlantedClique,Kearns1998SQ} for planted
problems and statistical queries.

\begin{theorem}[Seed catalogue size]\label{thm:calibration-family-structure}
The seed catalogue \(\mathcal C_{0}\) is a list of thirteen calibration
records; in particular \(|\mathcal C_{0}|\ge 13\).
\end{theorem}

\begin{proof}
The list in \Cref{tab:calibration:seed} has thirteen entries, each with the three fields of \Cref{def:calibration}; the last entry is a
transport between calibrations rather than a detector of its own.
\end{proof}

The catalogue is a seed, not a closed list: the record type is open,
and further calibrations can be adjoined, subject to the rebadge
criterion below.  We call the family \emph{non-descending} in the
intended sense that no entry is presented as a quotient or
specialization of another; this is a description of the entries, not a
proved property.

\begin{table}[H]
  \centering
  \caption{The seed catalogue of \Cref{thm:calibration-family-structure}.}
  \label{tab:calibration:seed}
  \scriptsize
  \renewcommand{\arraystretch}{1.06}
  \begin{tabularx}{\textwidth}{@{}>{\raggedright\arraybackslash}p{0.10\textwidth}>{\raggedright\arraybackslash}p{0.26\textwidth}>{\raggedright\arraybackslash}X>{\raggedright\arraybackslash}p{0.26\textwidth}@{}}
    \toprule
    Symbol & Name & Native carrier type & Detector class \\
    \midrule
    \(\Cvobs\) & volumetric observables & \(AC^0\) / \(AC^0[p]\)-adjacent carriers & low-degree polynomial detector \\
    \(\CMI\) & information calibration & bounded-information packages & mutual information / KL divergence \\
    \(\CRE\) & stability calibration & random-restriction / coupling-stable carriers & stable, local, overlap algorithms \\
    \(\CCh\) & cohomological calibration & cohomology-lifted carriers with Lov\'asz--Kneser reach & topological obstruction \\
    \(\CQamp\) & amplitude calibration & black-box and query carriers & quantum amplitude / query detector \\
    \(\CIalg\) & algebraic calibration & proof-system and certificate-degree carriers & ideal / Nullstellensatz detector \\
    \(\CRreal\) & realizability calibration & lambda-calculus / typed-program carriers & realizability detector \\
    \(\COCstar\) & \(C^*\)-operator calibration & operator-algebraic carriers & \(C^*\)-operator detector \\
    \(\CFexistsSO\) & existential-second-order calibration & logical / locality carriers & descriptive-complexity detector \\
    \(\CSHO\) & higher-order strategy calibration & higher-order game carriers & game-semantic strategy detector \\
    \(\CPPH\) & persistent-homology calibration & scale-dependent topological carriers & persistent-homology detector \\
    \(\CAVNP\) & arithmetic-complexity calibration & GCT / polynomial-family carriers & arithmetic-complexity / VNP detector \\
    \(\CMmot\) & computational motive (META) & cross-calibration transport carriers & motive-level transport detector \\
    \bottomrule
  \end{tabularx}
\end{table}


\subsection{Rebadges}\label{sec:calibration:rebadge}

A new name does not make a new calibration.  A candidate must differ
from every existing entry in the structural data by which a calibration
earns its place.

\begin{definition}[Calibration four-axis profile]\label{def:four-axis-profile}
The \emph{\fourAxis} of a calibration is the quadruple of labels
\[
  p=(p_{\mathsf{alg}},p_{\mathsf{det}},p_{\mathsf{desc}},p_{\mathsf{orient}})
\]
recording its algebraic identities, detector class, descent invariants
and carrier orientation.  A profile \(p\) is a \emph{rebadge} of \(q\),
\(\mathsf{Rebadge}(p,q)\), when \(p=q\) componentwise, and \(p\) is
\emph{genuinely new} relative to a family \(\mathcal F\),
\(\mathsf{GenuineNew}(p,\mathcal F)\), when it is a rebadge of no member
of \(\mathcal F\).
\end{definition}

\begin{theorem}[Rebadge Discipline]\label{thm:rebadge-discipline}
A profile \(p\) is genuinely new relative to \(\mathcal F\) if and only
if, for every \(q\in\mathcal F\), \(p\) and \(q\) differ on at least one
of the four axes.
\end{theorem}

\begin{proof}
\(\mathsf{Rebadge}(p,q)\) is the conjunction of the four componentwise
equalities, so its negation is the disjunction of the four
inequalities; quantify over \(q\in\mathcal F\).
\end{proof}

The criterion compares labels.  Whether two detectors with different
labels are mathematically equivalent is a separate question, which the
criterion does not decide.

\subsection{Bridges between calibrations}\label{sec:calibration:bridges-as-theorems}

\begin{definition}[Calibration-side typed bridge]\label{def:calibration-side-typed-bridge}
A \emph{\TypedBridgeCalibration} is a record
\(B=\{\mathsf{source},\mathsf{target},\mathsf{kind}\}\) consisting of a
source calibration, a target calibration and a kind.\footnote{This is
distinct from the carrier-side bridge of
\Cref{def:carrier-side-typed-bridge}.}
\end{definition}

\begin{definition}[Calibration-bridge kind dichotomy]\label{def:bridge-kind}
The kind of a calibration bridge is one of two values,
\(\mathsf{BridgeKind}::=\mathsf{basisChange}\mid\mathsf{substantive}\):
a \emph{basis change} relates two presentations of the same detector and
carrier content, and a \emph{substantive} bridge asserts a theorem
relating different invariants.
\end{definition}

\begin{theorem}[Bridge kinds]\label{thm:bridges-as-theorems}
Every calibration bridge record has kind \(\mathsf{basisChange}\) or
kind \(\mathsf{substantive}\).
\end{theorem}

\begin{proof}
\(\mathsf{BridgeKind}\) has exactly these two values.
\end{proof}

The content of the theorem is a convention: a transport between
calibrations is recorded only together with a declared kind, and is
then either a change of basis or a theorem that has to be proved.
Placing two calibrations side by side without such a declaration does
not produce a bridge record.  The kind label itself does not supply
the change of basis or the proof.

\subsection{Direction}\label{sec:calibration:bridge-type-direction}

Transports between calibrations usually have a direction: worst case to
average case, nonuniform to uniform, statistical to algorithmic.

\begin{definition}[Calibration-bridge reversal]\label{def:calibration-bridge-reversal}
The \emph{reversal} of a bridge record \(B=(S,T,k)\) is
\(B^{\mathrm{op}}=(T,S,k)\).
\end{definition}

\begin{theorem}[Bridge Type-Direction]\label{thm:bridge-type-direction}
For every calibration bridge record \(B\),
\(B=B^{\mathrm{op}}\) if and only if
\(\mathsf{source}(B)=\mathsf{target}(B)\).
\end{theorem}

\begin{proof}
If \(B=B^{\mathrm{op}}\), comparing source fields gives \(S=T\);
conversely, if \(S=T\) the two records coincide.
\end{proof}

Thus a bridge record between different calibrations is never its own
reversal.  Whether a bridge \(S\to T\) yields a usable bridge
\(T\to S\) is a mathematical question about the transport, which the
records do not model.

\subsection{Coverage}\label{sec:calibration:coverage-family-geometry}

Each calibration \(c\) is assigned a set \(\operatorname{Cov}(c)\) of
region labels, and the family coverage is
\(\operatorname{Cov}(\mathcal C_{0})=\bigcup_{c\in\mathcal C_{0}}\operatorname{Cov}(c)\).
The labels and the assignment form a finite declared table, shown in
\Cref{tab:coverage-regions}; the assignment is a modelling choice, and
the seven calibrations not listed there are assigned no region.  Two uncovered labels are singled out as
\emph{typed coverage holes}: \(\xiNCone\), with region
\(\mathsf{NC}^{1}\), and \(\xiAdaptive\), with region
\(\mathsf{adaptivePolyTime}\).  The standard references for the classes
suggested by the uncovered labels are \cite{Cook1985NCTaxonomy} for the
\(\mathsf{NC}\) hierarchy and \cite{KarpLipton1982Advice} for
\(\mathsf{P/poly}\).

\begin{table}[H]
  \centering
  \caption{The declared coverage table of \Cref{thm:coverage-family-geometry}.  The assignment of regions to calibrations is a modelling choice; calibrations not listed are assigned no region.}
  \label{tab:coverage-regions}
  \small
  \renewcommand{\arraystretch}{1.08}
  \begin{tabularx}{\textwidth}{@{}>{\raggedright\arraybackslash}p{0.30\textwidth}>{\raggedright\arraybackslash}X>{\raggedright\arraybackslash}p{0.22\textwidth}@{}}
    \toprule
    Region label & Assigned to & Status \\
    \midrule
    bounded-depth & \(\Cvobs\), \(\CMI\) & covered \\
    low-degree & \(\Cvobs\) & covered \\
    stable-local & \(\CRE\) & covered \\
    query & \(\CQamp\) & covered \\
    topological & \(\CCh\) & covered \\
    low-Nullstellensatz-degree & \(\CIalg\) & covered \\
    \(\mathsf{NC}^{1}\) & none & uncovered; hole \(\xiNCone\) \\
    \(\mathsf{NC}^{2}\) & none & uncovered \\
    \(\mathsf{P/poly}\) & none & uncovered \\
    adaptive polynomial time & none & uncovered; hole \(\xiAdaptive\) \\
    \bottomrule
  \end{tabularx}
\end{table}


\begin{theorem}[Declared coverage table]\label{thm:coverage-family-geometry}
In the declared table: the labels bounded-depth, low-degree,
stable-local, query, topological and low-Nullstellensatz-degree belong
to \(\operatorname{Cov}(\mathcal C_{0})\); the labels \(\mathsf{NC}^{1}\),
\(\mathsf{NC}^{2}\), \(\mathsf{P/poly}\) and adaptive polynomial time
do not; the regions of \(\xiNCone\) and \(\xiAdaptive\) are uncovered;
two distinct seed calibrations, \(\Cvobs\) and \(\CMI\), both cover the
bounded-depth label; and for every family \(\mathcal F\) and every
calibration \(c\) with \(\operatorname{Cov}(c)=\varnothing\),
\[
  \operatorname{Cov}(\mathcal F\cup\{c\})=\operatorname{Cov}(\mathcal F).
\]
\end{theorem}

\begin{proof}
The membership statements are checked entry by entry in the table.  The
last statement holds because the union with an empty set is unchanged.
\end{proof}

The theorem verifies the table; it proves no inclusion, exclusion or
separation between the complexity classes named by the labels.  Its
last clause is a reminder that adding names to the catalogue does not
by itself enlarge its coverage: a new calibration changes the coverage
only if it covers a new region, or if a bridge theorem carries content
into one.

\subsection{Native type and target type}
\label{sec:calibration:framework-type-self-char}

\begin{definition}[Framework-type]\label{def:framework-type}
A \emph{\FType} is a record of four labels: uniformity, worst-caseness,
scoping and packaging.  Two framework-types \emph{match} when they agree
on all four.  For the P-versus-NP application we record
\[
\begin{aligned}
  F_{\mathrm{PvNP}}
  &=(\mathsf{nonuniform},\mathsf{averageCase},\mathsf{scoped},\mathsf{packageBased}),\\
  T_{\mathrm{Clay}}
  &=(\mathsf{uniform},\mathsf{worstCase},\mathsf{unconditional},\mathsf{languageBased}).
\end{aligned}
\]
\end{definition}

The first two axes correspond to the distinctions surveyed in
\cite{Impagliazzo1995FiveWorlds}.  The tuple \(F_{\mathrm{PvNP}}\)
describes the calibration catalogue of this section, not every
construction of the application: the P-versus-NP paper itself proves a
uniform, worst-case machine translation and a conditional separation of
standard languages.  The comparison below concerns the two recorded
tuples and does not classify that proof.

\begin{theorem}[Recorded type mismatch]\label{thm:framework-type-self-char}
The recorded types \(F_{\mathrm{PvNP}}\) and \(T_{\mathrm{Clay}}\) do not
match, and the gate catalogue of \Cref{cat:named-functorial-gates} has
three entries.
\end{theorem}

\begin{proof}
The first labels, \(\mathsf{nonuniform}\) and \(\mathsf{uniform}\), are
different (so are the other three).  The count is
\Cref{cat:named-functorial-gates}.
\end{proof}

\begin{schema}[Framework-general candidate]\label{sch:framework-type-pending}
For every formed closure \(c\) in the trace-state-only column, the
native framework-type assigned to \(c\) and the framework-type of its
standard target do not match.
\end{schema}

Both type assignments in the schema are primitive.  The schema is not
used in the proof of \Cref{thm:framework-type-self-char}; it records the
conjectured extension beyond the P-versus-NP case.  The mismatch says
where the catalogue's native reach lies, and therefore which transports
would have to be supplied to reach the standard target; the next
subsection names three of them.

\subsection{Functorial gates}\label{sec:calibration:functorial-gates}

\begin{definition}[Named functorial gate]\label{def:functorial-gate}
A \emph{named functorial gate} is a record
\[
  \{\mathsf{name},\mathsf{source},\mathsf{target},\mathsf{direction},\mathsf{theoremObligation}\}
\]
of labels: a typed source, a typed target, a direction, and the
statement of the theorem that a transport from source to target would
have to prove.  The record does not assert that the transport exists.
\end{definition}

\begin{table}[H]
  \centering
  \caption{The three named functorial gates of \Cref{cat:named-functorial-gates}.  Each records a transport that has not been proved.}
  \label{tab:functorial-gates}
  \small
  \renewcommand{\arraystretch}{1.1}
  \begin{tabularx}{\textwidth}{@{}>{\raggedright\arraybackslash}p{0.08\textwidth}>{\raggedright\arraybackslash}p{0.22\textwidth}>{\raggedright\arraybackslash}p{0.22\textwidth}>{\raggedright\arraybackslash}X@{}}
    \toprule
    Gate & Source & Target & Obligation \\
    \midrule
    \(\GateFunctionDegree\) & degree of a low-degree detector (\(\Cvobs\)) & degree of a Nullstellensatz certificate & relate the degree seen by the detector to the degree of an ideal-membership certificate \\
    \(\GatePlantedStatistical\) & average-case (planted) statistical detection & worst-case hardness of an NP-hard language & carry detection content from the average-case side to the worst case; as a hardness statement, show that the language being hard in the worst case makes the detection problem hard on average \\
    \(\GateUniformTime\) & nonuniform families & uniform time hierarchy & realize a nonuniform family in a uniform time-hierarchy argument of the kind used by Williams \cite{Williams2014NonUniformACC} \\
    \bottomrule
  \end{tabularx}
\end{table}


\begin{theorem}[Gate catalogue]\label{cat:named-functorial-gates}
The gate catalogue \(\mathcal G_{\mathrm{named}}=[G_{1},G_{2},G_{3}]\)
of \Cref{tab:functorial-gates} has three entries, and each has a
nonempty name and a nonempty theorem obligation.
\end{theorem}

\begin{proof}
Inspection of the three records.
\end{proof}

None of the three obligations is proved here; they are the open
problems of \Cref{sec:predictions:functorial-gates}.  The gate records
complement the classical barrier results---relativization, natural
proofs and algebrization
\cite{BakerGillSolovay1975Relativization,RazborovRudich1997NaturalProofs,AaronsonWigderson2009Algebrization}---which
show that certain families of techniques cannot cross certain
boundaries; a gate states what a crossing would have to prove.  The
P-versus-NP application also contains conditional results restricting
feature-based proof attempts, under explicit hypotheses: a
feature-decomposition record, a joint-access operation, and assumptions
that expressions of given classes do not prove the target
\cite{Tsiokos2026PvNP}.  These results do not constitute a barrier of
the scope of the classical ones.  The three gates are examples, not a complete list.

\section{Cross-paper consistency}\label{sec:audit}

This section adds no formal claim.  It checks that the way this paper
describes the companion papers agrees with what those papers state,
and that the two recognition sources of \Cref{sec:recognition} are
recorded as named premises rather than re-derived.

\subsection{Consistency of the framing}\label{sec:audit:cross-paper-consistency}

\paragraph{Recognition and derivation.}  In this paper a
recognition-mode landing and a derivation-mode landing are two
different kinds of conditional closure: the first names its
load-bearing source as a premise, the second derives its conclusion
from stated hypotheses.  Neither is treated as a weaker form of the
other.  This is consistent with the companion papers.  The
Riemann-hypothesis paper \cite{Tsiokos2026RH} proves the Riemann
hypothesis from a supplied Selberg recognition source and an
identification with the classical zeros; the P-versus-NP paper
\cite{Tsiokos2026PvNP} proves a fixed machine-level translation and
derives \(P\neq NP\) from an accepted negative source.  Both state these
premises as hypotheses.  The same separation is what keeps
\Cref{thm:attack-foreclosure-v5} from reading as a retraction of
\Cref{thm:attack-foreclosure-v4}: the first applies the second to
derivation routes only.

\paragraph{Column placement.}  The Riemann-hypothesis and P-versus-NP
substrates are placed in the trace-state-only column, and the
Navier--Stokes substrate in the witness-content-exposing column, by the
V-Differential classification of the hiddenness companion
\cite{Tsiokos2026CSL}.  That paper adopts the placements as
hypotheses, with an informal motivation for each; it does not prove
them.  This paper uses the same placements with the same status.  The
two trace-state-only substrates are mathematically unlike each other,
one analytic and one computational, but the criterion applied to them
is the same.

\paragraph{Source handling.}  In the P-versus-NP row the named source
is
recorded in the P-versus-NP paper \cite{Tsiokos2026PvNP}; in the
Riemann-hypothesis row the Selberg source is recorded in the
Riemann-hypothesis paper \cite{Tsiokos2026RH}.  In both rows the source is a named
premise and is not produced by a derivation in this paper.  A source
record is broader than the proposition it contributes: it carries its
provenance, the closure it belongs to and its column placement, while
the proposition is what enters the conditional argument.

\subsection{Recording sources}\label{sec:audit:source-record-discipline}

We use the following rule.  A recognition source must be named, its
substrate must be identified, and its role in a closure must be kept
separate from the proposition it contributes.  A proposed
recognition-mode landing that cannot name its source, identify its
column placement, or separate the source from the target proposition is
not counted as an instance of \Cref{thm:template-applicability}; such
cases belong to the open problems of \Cref{sec:predictions}.

For the P-versus-NP row, naming the source is bookkeeping, not
evidence for its negative content.  The hiddenness companion proves a
finite example in which two histories have different predictive classes
and the same current class while selected Boolean outputs read from the
two quotients agree; it does not connect its quotients with the classes
\(P\) and \(NP\).  The accepted negative content and the admission of
search events used for the standard complexity conclusion are separate
premises of the P-versus-NP paper.

\section{Predictions and open problems}\label{sec:predictions}

This section records predictions and open problems.  None of them is
used as evidence for a numbered statement, and none asserts that a
missing application has been carried out.

\subsection{A Hodge substrate}
\label{sec:predictions:hodge}

We conjecture, on structural grounds only, that a Hodge substrate
\cite{Deligne2000Hodge,Voisin2003HodgeTheoryII} would lie in the
trace-state-only column; no formed closure for the Hodge conjecture has
been constructed in this series.  If one is constructed and the column
placement holds, the seven-stage template of \Cref{sec:recognition}
should apply to it, and its derivation sweep would be a third test of
the verdict signature \SigPattern\ of \Cref{thm:verdict-signature}.
The prediction concerns the column, not the substrate: it asserts no
Hodge closure and names no Hodge recognition source.  It would be
refuted by a trace-state-only Hodge closure satisfying the template
whose derivation sweep returns a different verdict pattern; if a Hodge
substrate turned out not to be trace-state-only, the prediction would
simply not apply.

\subsection{Yang--Mills}
\label{sec:predictions:ym}

Yang--Mills \cite{JaffeWitten2000YangMills} is a further candidate.
No Yang--Mills substrate has been studied in this series, and nothing
is claimed about it.  The open question is whether a Yang--Mills
substrate would be trace-state-only, witness-content-exposing, or of a
kind not yet represented.  The column question should be settled before
a recognition-mode or derivation-mode template is chosen.

\subsection{Proving the recognition-mode schemas}
\label{sec:predictions:theorem-grade-lifting}

The schemas of \Cref{sec:recognition} can become theorems only after
their primitive predicates are defined; until then they are neither
provable nor refutable.  Given definitions, each schema needs its own
argument.  \Cref{thm:template-applicability} needs a proof that a
successful seven-stage audit gives the closure the status of a
conditional theorem under its premises.  \Cref{thm:verdict-signature}
needs a derivation of the verdict pattern \SigPattern\ from column
membership.  \Cref{thm:closure-content-thesis} needs, for the
load-bearing content of a trace-state-only closure, a proof that it is
formed-layer content and is named as a source, and a non-derivability
proof showing that the framework's primitives do not yield it.
\Cref{thm:attack-foreclosure-v5} needs, beyond
\Cref{thm:attack-foreclosure-v4}, a proof that every recognition-mode
landing imports the load-bearing content and lands at theorem grade.
Further examples, such as a third trace-state-only application beyond
the P-versus-NP and Riemann-hypothesis cases, would test these
arguments but cannot replace them.  This is the main open problem left
by \Cref{sec:recognition}.

\subsection{A bound on the depth of matrix-element towers}
\label{sec:predictions:ctmt-depth-bound}

\Cref{thm:ctmt-recursion} describes towers of matrix-element
terminations of any finite depth but gives no bound on the depth.  The
open problem is to find an invariant of the carrier
that bounds the depth independently of the application, for example the
complexity of the inherited records, a descent height, the rank of the
native matrix family, or a refinement of the adequacy residual.  Such a
bound cannot simply forbid further reductions by definition; it has to
explain why a deeper matrix-element family cannot arise in a typed
context.

\subsection{Open functorial gates}
\label{sec:predictions:functorial-gates}

The three gates of \Cref{cat:named-functorial-gates},
\(\GateFunctionDegree\), \(\GatePlantedStatistical\) and
\(\GateUniformTime\), each name a transport that has not been proved: from the degree of a
detector to the degree of a Nullstellensatz certificate, from planted
statistical detection to worst-case hardness of a language, and from a
nonuniform family to a uniform time-hierarchy statement.  Each is a
problem in its own right.  The second transports content from the
average case to the worst case, so the corresponding hardness statement
is a worst-case to average-case reduction, and it is related to known
obstacles: a
non-adaptive reduction from an NP-complete problem in the worst case to
a distributional NP problem under a polynomial-time samplable
distribution, tolerant to a small fraction of oracle errors, would
imply \(\mathrm{coNP}\subseteq\mathrm{NP/poly}\)
\cite{BogdanovTrevisan2006WorstAverage}, extending the earlier result
of \cite{FeigenbaumFortnow1993RandomSelfReducibility}; whether a
planted detection problem falls under these hypotheses depends on its
formulation.  Further gates should appear whenever another
application needs a transport between calibrations, regimes, or a
native type and a target type that is not a change of basis, for
instance when its coverage table has a hole like those of
\Cref{thm:coverage-family-geometry}.  The list of gates is not claimed
to be complete.

\section{Scope and limitations}\label{sec:scope}

\subsection{Summary}\label{sec:scope:nonclaims}

This section collects what the paper does not claim.
\Cref{tab:nonclaims} summarizes it, and the subsections below give the
details; they govern how the earlier sections should be read.

\begin{table}[H]
  \centering
  \caption{What the paper does not claim.}
  \label{tab:nonclaims}
  \small
  \renewcommand{\arraystretch}{1.1}
  \begin{tabularx}{\textwidth}{@{}>{\raggedright\arraybackslash}p{0.30\textwidth}>{\raggedright\arraybackslash}X@{}}
    \toprule
    Topic & Limitation \\
    \midrule
    Birch--Swinnerton-Dyer & Outside the scope of this paper; not used as evidence. \\
    Hodge & A prediction only (\Cref{sec:predictions:hodge}). \\
    Unconditional results & None.  Each application keeps its own premises; the Gaussian constructions supply none of them. \\
    Recognition sources & Named premises, not derived; their acceptance and any admission laws are hypotheses. \\
    Recognition-mode schemas & Conditional schemas illustrated by two applications; not proved. \\
    Non-descending translation & Not stated or used as a theorem; represented only by the calibration family and the open gates. \\
    \bottomrule
  \end{tabularx}
\end{table}


\subsection{The Birch--Swinnerton-Dyer conjecture}\label{sec:scope:bsd-out}

The paper treats three applications: Navier--Stokes regularity, the
Riemann hypothesis, and \(P\) versus \(NP\).  A companion paper treats
the Birch--Swinnerton-Dyer conjecture \cite{TsiokosBSD2026,Wiles2000BSD};
it is not analysed here and is not used as evidence.  This is a choice
of scope, not a statement about that paper.

\subsection{The Hodge conjecture}\label{sec:scope:hodge-prediction-only}

No formed closure for the Hodge conjecture exists in this series.  The
Hodge conjecture appears only in the prediction of
\Cref{sec:predictions:hodge}, which proposes it as a third test case.
The paper does not claim that the conjecture holds, that it satisfies
the recognition-mode template, or that a Hodge recognition source has
been found.

\subsection{No unconditional claim}\label{sec:scope:standard-zfc}

Every application cited here is conditional on its own premises, and
the framework account also uses the SB closure assumption.  This paper
strengthens none of the applications to an unconditional theorem and
does not describe any of them as unconditional.  The premises are, for
Navier--Stokes, a spectral closure, or a physical-window closure
together with an assumed terminal-divergence input, with initial-data
hypotheses; for the Riemann hypothesis, a supplied
Selberg recognition source and an identification with the classical
zeros; and for \(P\) versus \(NP\), that a fixed tagged witness
selector is not computable in polynomial time, or, in the family form,
a named negative source together with the admission of search events.

\Cref{thm:joint-gaussian-return} is unconditional as an identity on the
objects it names.  Its zero-side premise \(\mathcal A_{a,\nu,t}\) is
equivalent to the Riemann hypothesis, and its Navier--Stokes conclusions
assume a separate closure hypothesis; the identity establishes neither
and does not combine them into one closure law.
\Cref{thm:concrete-triad-gaussian-xi} likewise proves exact identities.
Its SAT sum has \(2^{n}\) terms, and its positivity gives no
polynomial-time procedure.  It transfers no premise from one
application to another, and the impossibility statement of
\Cref{sec:triad-bridge} excludes only maps that preserve the probe and
send the reflected displacement to a satisfiable SAT state.

\subsection{Recognition sources are premises}\label{sec:scope:recognition-source-imports}

The recognition sources are named premises, not derived results.  The
P-versus-NP source is recorded in that paper; its acceptance and the
admission of search events are separate premises.  The Selberg source
is recorded in the Riemann-hypothesis paper.
The paper never treats these sources as derived from the framework or
produced by the recognition-mode template.

\subsection{The recognition-mode schemas are not proved}\label{sec:scope:universal-lifting-deferred}

The four statements of \Cref{sec:recognition} are conditional schemas,
illustrated by the two companion applications and not proved.
Turning them into theorems needs definitions of their predicates and
the arguments listed in \Cref{sec:predictions:theorem-grade-lifting}.  The relation between
\Cref{thm:attack-foreclosure-v5} and \Cref{thm:attack-foreclosure-v4}
is unchanged by this: the former restricts the latter to derivation
routes and adds the recognition-mode clauses.

\subsection{Non-descending translation}
\label{sec:scope:m6-theorem-7-folded}

No non-descending translation theorem for the P-versus-NP application
is stated or used here.  The transport content such a statement would
carry is represented in this paper only by the calibration family of
\Cref{sec:calibration} and the open gates of
\Cref{sec:predictions:functorial-gates}.


\appendix
\section{Formal verification}\label{app:formalization}

\subsection{Setting}\label{app:formalization:setting}

The results of this paper are formalized in Lean~4
\cite{deMouraUllrich2021Lean4}, version 4.28.0, with Mathlib
\cite{MathlibCommunity2020Lean} at revision \texttt{8f9d9cff}.  The
Gaussian results of \Cref{sec:gaussian-bridge,sec:triad-bridge} use the
Lean libraries of the companion papers: the Navier--Stokes and
adequacy-residual libraries for mild solutions, the heat semigroup and
the residual projection \(\Xi(L)\), the Riemann-hypothesis library for
the zeros of the completed zeta function and their reflection, and the
P-versus-NP library, built on complexitylib \cite{Schlesinger2026Complexitylib},
for CNF semantics and the bitstring encoding of \(L_{\mathrm{SAT}}\).
The code is in the repositories
\url{https://github.com/ioannist/six-birds-meta-math},
\url{https://github.com/ioannist/six-birds-needles},
\url{https://github.com/ioannist/six-birds-duality-confinement} and
\url{https://github.com/ioannist/six-birds-hiddenness}; the first
depends on the other three, which are expected to be checked out side
by side with it.  As in the Navier--Stokes companion,
velocity fields are represented by their Sobolev-weighted Fourier
states (\Cref{sec:gaussian-bridge}), and the mild solution and the
Duhamel term are defined in these coordinates.  The zeros of the
completed zeta function, the completed zeta function itself and the
statement of the Riemann hypothesis are those of Mathlib.

The formal counterparts of \Cref{sec:foreclosure,sec:recognition} are
written as follows.  Claims are elements of a type with an
interpretation into propositions, so that distinct claims can have the
same truth value (\Cref{sec:apparatus:closure-formation}).  Primitive
predicates are either parameters of the statements, collected in
signature records, or opaque constants about which nothing is assumed.
Each conditional schema is a definition of a proposition; it is never
asserted, and the results that use a schema take it as an explicit
hypothesis.  For the Coverage schema the formal definition is the
per-instance conclusion, without the aiming, typing and coverage-context
guards of the displayed statement; it is thus a stronger proposition,
used only as a hypothesis.  Primitive
predicates of \Cref{sec:recognition}, including the template-applies
predicate, are opaque constants; the seven stage propositions of a
template record are not used in their definition.  Formal verification
therefore confirms the logical shape of the statements of
\Cref{sec:foreclosure,sec:recognition,sec:calibration} and the
identities of \Cref{sec:gaussian-bridge,sec:triad-bridge}; it confirms
no interpretation of the framework's predicates.  In \Cref{sec:foreclosure} the carrier residual is a matrix
over the rationals given by an explicit formula; only the proposition
\(R_{C}=0\) is used.

\subsection{Correspondence}\label{app:formalization:master-mapping}

\Cref{tab:formal-correspondence} lists the formal counterpart of every
numbered statement.  Identifiers are relative to the namespace
\leanid{SixBirdsMetaMath}; we abbreviate
\leanid{MetaMath.Foreclosure} as \leanid{F},
\leanid{MetaMath.RecognitionMode} as \leanid{R} and
\leanid{MetaMath.CalibrationFamily} as \leanid{C}.  The statuses are:
\emph{proved}, a theorem with no hypotheses beyond those stated in the
paper; \emph{proved from} (Cls), (BO) or coverage, a theorem taking that
schema instance as a hypothesis; \emph{schema}, a definition of the
schema as a proposition, not asserted; \emph{definition} and
\emph{primitive}, for definitions and for primitive predicates.

\begingroup\small
\captionsetup{justification=raggedright,singlelinecheck=false}
\begin{longtable}{@{}>{\raggedright\arraybackslash}p{0.16\textwidth}>{\raggedright\arraybackslash}p{0.54\textwidth}>{\raggedright\arraybackslash}p{0.24\textwidth}@{}}
\caption{Formal counterparts of the numbered statements.}\label{tab:formal-correspondence}\\
\toprule
Statement & Declarations & Status\\
\midrule
\endfirsthead
\toprule
Statement & Declarations & Status\\
\midrule
\endhead
\Cref{def:carrier-residual} & \leanid{F.CarrierDichotomy.carrierResidual} & definition (matrix form)\\
\Cref{def:typed-primary-carrier} & \leanid{F.CarrierDichotomy.ClaimLanguage}, \leanid{F.CarrierDichotomy.TypedPrimaryCarrier} & definition\\
\Cref{def:four-typed-states} & \leanid{F.CarrierDichotomy.xEquivalent}, \leanid{nonXEquivalent}, \leanid{creTE}, \leanid{creBF}, \leanid{creCTMT}, \leanid{NativeClosureOutcome} & definition\\
\Cref{def:coverage-ctx} & \leanid{F.CarrierDichotomy.CoverageContext} & primitive\\
\Cref{thm:carrier-dichotomy} & \leanid{F.CarrierDichotomy.carrier_dichotomy}; (Cls) is \leanid{CarrierClassificationSchema} & proved from (Cls)\\
\Cref{def:carrier-side-typed-bridge} & \leanid{F.BridgeImpossibility.BridgeSignature}, \leanid{TypedBridge}, \leanid{BridgeComposition} & primitive\\
\Cref{def:x-strength-obligation} & \leanid{F.BridgeImpossibility.XStrengthObligation} & primitive\\
\Cref{cor:bridge-impossibility} & \leanid{F.BridgeImpossibility.bridge_impossibility}; (BO) is \leanid{BridgeObligationSchema}; joint satisfiability: \leanid{bridge_impossibility_hypotheses_satisfiable} & proved from (Cls) and (BO)\\
\Cref{thm:coverage} & \leanid{F.Coverage.CoverageSchema}; derivation from (Cls): \leanid{F.Coverage.coverage_of_classification} & schema\\
\Cref{thm:type-finite-foreclosure} & \leanid{F.TypeFiniteForeclosure.type_finite_foreclosure}; the choice of state: \leanid{classify}, \leanid{classify_spec} & proved from coverage\\
\Cref{def:in-scope} & \leanid{F.AttackForeclosure.InternalMoveSignature}, \leanid{InFrameworkScope}, \leanid{advancesXWithoutExternal} & primitive\\
\Cref{thm:attack-foreclosure-v4} & \leanid{F.AttackForeclosure.AttackForeclosureV4Schema} & schema\\
\Cref{thm:ctmt-recursion} & \leanid{F.CTMTRecursion.ctmt_recursion} & proved\\
\Cref{def:vdiff-tso-column} & \leanid{R.Template.VDifferentialTSOColumn} & primitive; the second column is not formalized\\
\Cref{def:applies} & \leanid{R.Template.SevenStageTemplate}, \leanid{TemplateApplies} & record; primitive\\
\Cref{thm:template-applicability} & \leanid{R.Template.TemplateApplicabilitySchema} & schema\\
\Cref{def:three-option-sweep} & \leanid{R.VerdictSignature.SweepOption}, \leanid{Verdict}, \leanid{sweepVerdict} & definition; primitive map\\
\Cref{thm:verdict-signature} & \leanid{R.VerdictSignature.VerdictSignatureSchema} & schema\\
\Cref{def:load-bearing} & \leanid{R.ClosureContentThesis.loadBearingFor} & primitive\\
\Cref{thm:closure-content-thesis} & \leanid{R.ClosureContentThesis.ClosureContentThesisSchema} & schema\\
\Cref{thm:attack-foreclosure-v5} & \leanid{R.AttackForeclosureV5.AttackForeclosureV5Schema}; first clause from \Cref{thm:attack-foreclosure-v4}: \leanid{attack_v5_derivation_conjunct_of_v4} & schema\\
\Cref{thm:joint-gaussian-return} & \leanid{RHNavier.JointSixBirdsObservationReturn.joint_actual_zero_xi_and_navier_source_return}; parts (1) and (2) without the fluid hypotheses: \leanid{RHNavier.ComplexGaussianNavierRHBridge.navierHeatWindowExcess_eq_rh_xi}, \leanid{actual_zero_gaussian_eq_navier_heat_readout} & proved\\
\Cref{thm:concrete-triad-gaussian-xi} & \leanid{RHNavierPvNP.TriadObservationReturn.concrete_triad_gaussian_xi_span}; normalized frame value: \leanid{RHNavier.ComplexGaussianNavierRHBridge.normalizedNonlinearFrameReadout}; existence of a positive admissible frame time: \leanid{RHNavierPvNP.ExternalSATGaussianLanguage.exists_positive_common_theta_time} & proved\\
\Cref{def:calibration} & \leanid{C.Structure.Calibration}, \leanid{calibrationFamily} & definition\\
\Cref{thm:calibration-family-structure} & \leanid{C.Structure.calibration_family_structure} & proved\\
\Cref{def:four-axis-profile} & \leanid{C.RebadgeDiscipline.CalibrationProfile}, \leanid{isRebadgeOf}, \leanid{genuinelyNew} & definition\\
\Cref{thm:rebadge-discipline} & \leanid{C.RebadgeDiscipline.rebadge_discipline} & proved\\
\Cref{def:calibration-side-typed-bridge} & \leanid{C.BridgesAsTheorems.TypedBridge} & definition\\
\Cref{def:bridge-kind} & \leanid{C.BridgesAsTheorems.BridgeKind} & definition\\
\Cref{thm:bridges-as-theorems} & \leanid{C.BridgesAsTheorems.bridges_as_theorems} & proved\\
\Cref{def:calibration-bridge-reversal} & \leanid{C.BridgeTypeDirection.TypedBridge.reverse} & definition\\
\Cref{thm:bridge-type-direction} & \leanid{C.BridgeTypeDirection.bridge_type_direction} & proved\\
\Cref{thm:coverage-family-geometry} & \leanid{C.CoverageFamilyGeometry.coverage_family_geometry}, \leanid{all_covered_regions_in_family}, \leanid{all_uncovered_regions_not_in_family}, \leanid{bounded_depth_overlap}, \leanid{empty_coverage_adds_nothing} & proved\\
\Cref{def:framework-type} & \leanid{C.FrameworkTypeSelfChar.FrameworkType}, \leanid{typesMatch}, \leanid{pvnpFrameworkType}, \leanid{standardClayPvNPType} & definition\\
\Cref{thm:framework-type-self-char} & \leanid{C.FrameworkTypeSelfChar.framework_type_self_characterization} & proved\\
\Cref{sch:framework-type-pending} & \leanid{C.FrameworkTypeSelfChar.FrameworkTypeGeneralSchema} & schema\\
\Cref{def:functorial-gate} & \leanid{C.FunctorialGates.FunctorialGate} & definition\\
\Cref{cat:named-functorial-gates} & \leanid{C.FunctorialGates.named_functorial_gates} & proved\\
\bottomrule
\end{longtable}
\endgroup

Some unnumbered statements are also formalized.  In
\Cref{sec:apparatus,sec:foreclosure}: a sound set of accepted claims
that accepts one true claim and rejects another
(\leanid{F.CarrierDichotomy.sound_base_noncollapse}), and a carrier that
is both \nonXequiv\ and \Xequiv\
(\leanid{F.CarrierDichotomy.nonXEquivalent_and_xEquivalent}).  In
\Cref{sec:gaussian-bridge}: the equivalence of \(\mathcal A_{a,\nu,t}\)
with the Riemann hypothesis
(\leanid{RHNavier.JointSixBirdsObservationReturn.arithmetic_zero_window_closure_iff_rh})
and the two conditional Navier--Stokes returns
(\leanid{navier_global_mild_of_two_channel_closure},
\leanid{navier_pointwise_physical_pde_of_hilbert_xi_closure}).  In
\Cref{sec:triad-bridge}: the impossibility of a probe-preserving map
(\leanid{RHNavierPvNP.ProbePreservingTransportNoGo.no_probe_preserving_transport_to_sat})
and the identification of the formulas with positive Gaussian mass with
\(L_{\mathrm{SAT}}\)
(\leanid{RHNavierPvNP.ExternalSATGaussianLanguage.gaussianSATLanguage_eq_L_SAT}).

\subsection{What is not formally verified}\label{app:formalization:not-verified}

\begin{itemize}
\item The conditional schemas.  They are stated, not proved, and no
  result of the paper assumes one except as an explicit hypothesis.
\item The meaning of the primitive predicates.  A result about a
  primitive predicate holds for every interpretation of it; that an
  intended interpretation has the stated properties is not verified.
\item The column placements, which are hypotheses of the hiddenness
  companion, and the verdicts of \Cref{tab:verdict-signature}, which are
  this paper's reading of the applications.
\item The descriptions of the companion papers in
  \Cref{sec:validation,sec:audit}, and the interpretations of the
  calibration labels and coverage regions in \Cref{sec:calibration}.
\item The statements of classical results quoted in the text.
\item That the library definitions of weighted Fourier states, mild
  solutions, the Duhamel term, the heat semigroup and the residual
  projection agree with the classical objects they model; this is a
  matter of reading the definitions, which are those of the companion
  libraries.
\end{itemize}

\subsection{Trust base}\label{app:formalization:trust}

Every declaration named in this appendix depends only on Lean's
standard axioms \texttt{propext}, \texttt{Classical.choice} and
\texttt{Quot.sound}; the modules of this paper declare no further
axioms.  Beyond these axioms, the trusted base consists of the Lean
kernel and the libraries named above.  The repository also contains a
separate development for another paper of the series, which likewise
declares no axioms.

\section{Version history}\label{app:version-history}
\begingroup\small
Every numbered statement of earlier versions keeps its number and
place.  Several statements were renamed in Version~4, as recorded
below.

\subsection*{Version 4 (October 1, 2026)}
\begin{itemize}
    \item The paper was retitled; Versions~1--3 were titled \emph{One
    Meta-Theory, Three Clay-Problem Closures}.  The abstract, the
    introduction and the scope section were rewritten to state the
    status of every result, and a glossary was added.
  \item In \Cref{sec:foreclosure}, targets and accepted imports are now
    claims with an interpretation, instead of propositions.  With
    propositions, a sound set of accepted imports could only accept all
    true propositions or none, and the hypotheses of
    \Cref{cor:bridge-impossibility} could not hold together; the formal
    development now exhibits a model in which they do.
    \Cref{thm:carrier-dichotomy} no longer assumes the aiming, typing
    and coverage-context hypotheses, which its proof does not use.  The
    classification hypothesis (Cls) and the bridge-obligation hypothesis
    (BO), previously applied as unproved obligations, are now explicit
    hypotheses of \Cref{thm:carrier-dichotomy,cor:bridge-impossibility},
        and \Cref{thm:type-finite-foreclosure} assumes coverage explicitly.
    It is shown that the Coverage schema follows from (Cls) when some
    claim expresses the residual closure.  The names of these three
    results now end in ``conditional form''.
  \item The formal record no longer states any schema as an axiom; each
    schema is a definition, used only as a hypothesis.
  \item \Cref{thm:attack-foreclosure-v4,thm:attack-foreclosure-v5},
    previously called Attack Foreclosure v4 and v5, are now called
    Attack Foreclosure and Refined Attack Foreclosure.  Readings of them
    that do not follow from their statements are now described as
    interpretations, and the first clause of
    \Cref{thm:attack-foreclosure-v5} is derived from
    \Cref{thm:attack-foreclosure-v4}.
    \item Definitions~\ref{def:vdiff-tso-column}, \ref{def:applies} and
    \ref{def:load-bearing}, previously called V-Differential
    trace-state-only column, Template-applies predicate and
    Load-bearing predicate, are now called Columns, Template applies and
    Load-bearing content; \Cref{thm:concrete-triad-gaussian-xi},
    previously the Concrete Gaussian--XI triad span, is now the Concrete
        Gaussian--residual triad.  In \Cref{sec:calibration},
    Theorems~\ref{thm:calibration-family-structure},
    \ref{thm:bridges-as-theorems}, \ref{thm:coverage-family-geometry},
    \ref{thm:framework-type-self-char} and
    \ref{cat:named-functorial-gates}, previously called
    Calibration-Family Structure, Bridges-as-Theorems, Coverage-Family
    Geometry, Framework-Type Self-Characterization and Named Functorial
    Gates, are now called Seed catalogue size, Bridge kinds, Declared
    coverage table, Recorded type mismatch and Gate catalogue, names
    that describe what they check.  \Cref{def:load-bearing} defines the
    load-bearing content as an assignment, as in the formal record.  The verdicts of
    \Cref{tab:verdict-signature} are described as this paper's reading
    of the two applications rather than as verified instances.
    \item The SB closure assumption is no longer attributed to
    Foundations~I, the result quoted from Foundations~III is stated
    with its scope, and the name of the Selberg source is described as
    an analogy.
  \item The recognition sources are attributed to the applications
    that name them, the formed-closure predicate is identified as a
    definition of this paper, the sweep reports read in
    \Cref{sec:validation} are cited to the first versions of the
    applications, and the Navier--Stokes physical-window return is
    stated with its time-integrability condition and its assumed
    terminal-divergence input.
  \item \Cref{sec:gaussian-bridge,sec:triad-bridge} now define the
        observation, the heat excess and the frame explicitly, in the
    Sobolev-weighted Fourier coordinates in which they are formalized.  The
    reflection is \(J(\rho)=1-\bar\rho\) (previously misprinted as
    \(1-\rho\)).  \Cref{thm:joint-gaussian-return} is stated for
    nonnegative integer weights, and \Cref{thm:concrete-triad-gaussian-xi}
    for the specific three-solution frame and unit weights for which it
    is proved; the earlier statement for an arbitrary frame was too
    general.
  \item In \Cref{sec:calibration} the conclusions of
    \Cref{thm:calibration-family-structure,thm:bridges-as-theorems,thm:bridge-type-direction}
    are restricted to the properties of records that are proved; the
    seed catalogue is described as introduced in this paper, and its
    source column was removed.
    \item The account of the non-descending translation statement, an
    earlier formulation from which the calibration family developed,
    was reduced to a statement of scope in \Cref{sec:scope}.
  \item Companion papers are cited in their current versions, and the
    first versions separately where their content is meant.
  \item A data point on the depth of matrix-element towers, not
    supported by the cited papers, was removed from
    \Cref{sec:predictions}.
    \item \Cref{app:formalization} was rewritten as a correspondence
    table with a statement of what is and is not verified.
    Cross-references now name the kind of statement they point to.
\end{itemize}

\subsection*{Versions 2 and 3}
Versions~2 and~3, the latter dated September 29, 2026, restated the
premises of the three companion applications in their current form,
added the Gaussian constructions of \Cref{sec:gaussian-bridge} and
\Cref{sec:triad-bridge} with their figures, and described their formal
verification.

\subsection*{Version 1 (June 17, 2026)}
The first version presented the foreclosure, recognition-mode and
calibration-family parts.
\par\endgroup


\begingroup\raggedright
\bibliographystyle{plainnat}
\bibliography{references}
\endgroup

\end{document}
